
\dictentry{C Star Efficiency}{Combustion efficiency measured by characteristic velocity, comparing the achieved c-star to the theoretical maximum for a given propellant combination. A high c-star efficiency indicates complete combustion and effective propellant utilization within the chamber volume. Typical high-performance engines achieve c-star efficiencies above 95 percent, and low values may signal injector or chamber design deficiencies. Improving c-star efficiency is a primary goal during engine development and hot-fire testing campaigns.}{catRM}

\dictentry{Canard Control}{Forward-mounted control surfaces for missile maneuvering that provide pitch and yaw authority ahead of the center of gravity. Canard fins generate lift at the nose to initiate rapid direction changes, giving the missile high agility against maneuvering targets. This configuration is popular in short-range air-to-air missiles where fast response times and tight turning radii are critical. The canard surfaces must be carefully sized to avoid aerodynamic coupling with the tail fins and to maintain stable flight throughout the engagement envelope.}{catRM}

\dictentry{Carbon-Carbon Nozzle}{A nozzle using carbon-carbon composite material that retains strength and resists erosion at extreme temperatures above 3000 K. Carbon-carbon nozzles are used on solid rocket motors and upper-stage liquid engines where refractory metals would be too heavy or would erode too quickly. The material is manufactured by infiltrating a carbon fiber preform with resin and then carbonizing it to achieve a dense, high-strength structure. These nozzles offer an excellent combination of thermal shock resistance and low erosion rates for high-performance propulsion applications.}{catRM}

\dictentry{Case}{The pressure vessel containing the solid propellant that must withstand internal combustion pressures and external flight loads throughout the burn. Motor cases are typically manufactured from high-strength steel, titanium, or filament-wound composite materials to minimize weight while providing adequate structural margins. The case also serves as the primary load-bearing structure of the motor, transmitting thrust loads to the airframe and supporting the grain during handling, transport, and flight operations.}{catRM}

\dictentry{Catalytic Igniter}{An igniter using a catalyst for hypergolic or monopropellant ignition by decomposing propellant on a reactive surface to generate heat. The catalyst bed raises the propellant temperature above its autoignition point, producing a hot gas stream that initiates main chamber combustion. Catalytic igniters are simple and reliable, commonly used in monopropellant thrusters and small liquid engines. They require no electrical power or pyrotechnic cartridges, making them ideal for applications where simplicity and long-term storability are essential.}{catRM}

\dictentry{CEP}{Circular error probable, a measure of accuracy defined as the radius of a circle centered on the target within which 50 percent of rounds are expected to land. A smaller CEP indicates greater precision and is a primary metric for comparing the accuracy of missile systems and artillery. Modern precision-guided munitions achieve CEPs of a few meters, while unguided ballistic missiles may have CEPs of hundreds of meters. Mission planners use CEP data to determine the number of missiles required to achieve a desired probability of target destruction.}{catRM}

\dictentry{Chaff}{Radar-reflective material for decoying consisting of thin metallic strips or wires cut to lengths that resonate at enemy radar frequencies. When released in a cloud, chaff creates a large radar cross-section that confuses tracking radar and seduces radar-guided missiles away from the actual target. Chaff has been used since World War II and remains an essential self-protection countermeasure on military aircraft and ships. Advanced chaff variants include calibrated corner reflectors and active repeating decoys for enhanced effectiveness against modern radar systems.}{catRM}

\dictentry{Chamber Pressure}{The pressure inside the combustion chamber resulting from the balance between propellant mass flow rate, combustion gas generation, and nozzle flow. Higher chamber pressure generally improves specific impulse by enabling more complete combustion and higher nozzle expansion ratios, but it also increases structural demands on the chamber and feed system. Typical liquid engines operate between 10 and 30 MPa, while solid motors range from 5 to 25 MPa. Chamber pressure is one of the most critical design parameters, directly influencing thrust level, engine efficiency, and structural mass.}{catRM}

\dictentry{Characteristic Velocity}{A measure of combustion performance, denoted c-star, defined as the product of chamber pressure and throat area divided by mass flow rate. Characteristic velocity depends on propellant chemistry, combustion temperature, and molecular weight of the exhaust gases, and is independent of nozzle geometry. It serves as the primary metric for evaluating injector and combustion chamber efficiency during hot-fire testing. Deviations from predicted c-star values indicate combustion quality issues such as incomplete mixing, poor atomization, or excessive heat loss.}{catRM}

\dictentry{Check Valve}{A valve allowing flow in one direction only by using a spring-loaded poppet, ball, or flapper that seats under reverse flow. Check valves prevent dangerous backflow of propellant, pressurant, or combustion gases that could damage upstream components or cause mixing of incompatible fluids. They are essential safety components placed throughout rocket propellant feed and pressurization systems. Reliable check valve performance is critical during engine start and shutdown transients when pressure reversals are most likely to occur.}{catRM}

\dictentry{Chugging}{Low-frequency combustion instability characterized by pressure oscillations typically below 200 Hz, often accompanied by visible pulsation of the exhaust flame. Chugging is caused by coupling between the combustion process and the propellant feed system, where pressure disturbances modulate propellant flow into the chamber. Anti-chugging measures include increasing feed system stiffness, adding baffles, and increasing injector pressure drops. If left unchecked, chugging can cause destructive mechanical vibrations, propellant feed disruption, and potential engine failure during operation.}{catRM}

\dictentry{Cigarette Burning}{End burning of a solid grain in which the combustion front progresses axially from one end of the cylindrical grain to the other along the longitudinal axis. This simple grain geometry produces a regressive thrust profile as the burning surface decreases with the receding flame front over time. Cigarette burning is used in applications requiring long burn times and low thrust levels, such as sustainer motors and certain tactical missile applications. The burn rate is controlled by propellant formulation and is highly predictable, making this configuration ideal for missions requiring consistent and extended thrust duration.}{catRM}

\dictentry{Circularization Burn}{A burn to circularize an orbit performed at the apogee of an elliptical transfer orbit to raise the perigee to match the apogee altitude. Without this maneuver the vehicle would continue on an elliptical path, and the circularization burn adds velocity perpendicular to the radius vector to equalize orbital altitude. This is a standard procedure in satellite deployment and is typically performed by the upper stage or spacecraft propulsion. The timing and magnitude of the burn must be precisely calculated to achieve the desired circular orbit within tight tolerance requirements.}{catRM}

\dictentry{Coast Phase}{The unpowered phase of missile flight between burnout of the booster and reentry or terminal guidance initiation. During coast, the missile follows a ballistic arc under gravitational and inertial forces with no active propulsion. The coast phase duration depends on the trajectory design and is a critical period for midcourse guidance corrections and target tracking by defense systems. For ballistic missile defense, the coast phase represents the primary engagement window for exoatmospheric interceptors attempting to destroy the warhead before reentry.}{catRM}

\dictentry{Coaxial Injector}{An injector with concentric fuel and oxidizer streams where one propellant flows through an annular gap surrounding a central post of the other propellant. This configuration promotes shear-driven mixing and atomization at the interface between the two streams. Coaxial injectors are widely used in high-performance liquid engines such as the SSME and Vulcain due to their good atomization characteristics and resistance to combustion instability. The annular gap geometry provides uniform propellant distribution across the injector face and tolerates a wide range of operating conditions.}{catRM}

\dictentry{Cold Gas Thruster}{A thruster using compressed gas for attitude control that expels stored pressurized nitrogen, helium, or other inert gas through a small nozzle. Cold gas systems are the simplest form of spacecraft propulsion with no combustion, no ignition hardware, and minimal complexity, making them ideal for small satellites and initial attitude acquisition. Their specific impulse is low, typically 50 to 80 seconds, limiting their use to short-duration or low-delta-v applications. Cold gas thrusters provide reliable, contamination-free operation for precision attitude control where simplicity outweighs performance.}{catRM}

\dictentry{Combustion Chamber}{The chamber where propellant burns in a rocket engine, designed to contain high-pressure combustion and provide sufficient residence time for complete reaction. The chamber volume and length-to-diameter ratio are optimized to balance complete combustion against heat loss and weight. Chamber walls are protected by regenerative cooling, film cooling, or ablative liners depending on engine type and burn duration. The chamber geometry directly affects combustion efficiency, stability characteristics, and the thermal loads experienced by the injector face and chamber walls.}{catRM}

\dictentry{Combustion Efficiency}{How completely propellant is burned, expressed as the ratio of actual to theoretical characteristic velocity or combustion temperature. High combustion efficiency indicates thorough mixing and reaction of fuel and oxidizer within the available chamber volume. It is a primary design objective of injector and chamber geometry optimization during engine development. Losses in combustion efficiency typically stem from poor atomization, inadequate propellant mixing, or insufficient residence time in the chamber, all of which can be addressed through injector redesign.}{catRM}

\dictentry{Combustion Instability}{Pressure oscillations in the combustion chamber caused by coupling between unsteady heat release and acoustic modes of the chamber or feed system. Combustion instability can range from low-frequency chugging to high-frequency screaming, and if unchecked it can cause catastrophic engine failure through thermal and mechanical overload. Suppression techniques include baffles, acoustic cavities, injector pattern changes, and Helmholtz resonators. Identifying and resolving combustion instability is often the most challenging and time-consuming phase of rocket engine development.}{catRM}

\dictentry{Combustion Stability}{Stable combustion without pressure oscillations, meaning the steady-state burn proceeds at nominally constant pressure and thrust without dangerous oscillatory behavior. Achieving combustion stability requires careful matching of injector design, chamber geometry, and propellant properties. Stability is verified through extensive hot-fire testing and is a critical qualification criterion before flight certification of any rocket engine. The absence of instability allows the engine to operate at its design performance level without risk of structural damage or thrust fluctuations during the burn.}{catRM}

\dictentry{Command Guidance}{Guidance from an external command link where tracking radar or sensors on the ground compute intercept solutions and transmit steering commands to the missile in flight. The missile acts as an obedient follower, executing course corrections sent from the launch platform without carrying its own target-seeking sensors. Command guidance simplifies the missile design but requires continuous radar illumination and is vulnerable to jamming of the command link. Despite these limitations, command guidance remains widely used for its simplicity and the ability to upgrade engagement algorithms without modifying the missile.}{catRM}

\dictentry{Contact Fuse}{A fuse detonating on impact that uses a mechanical or piezoelectric sensor to detect target contact and initiate the warhead. Simple contact fuzes are reliable and inexpensive, making them suitable for rockets, artillery, and some missile warheads. More advanced variants include setback and inertia switches to arm only after launch, preventing accidental detonation during handling and storage. Contact fuzes are often combined with proximity or delay modes in multi-mode fuze systems to maximize effectiveness across different target types.}{catRM}

\dictentry{Conventional Warhead}{A non-nuclear explosive warhead containing high-explosive filler such as Composition B or PBX that generates blast overpressure and fragmentation to damage the target. Conventional warheads can be configured as blast, fragmentation, or combined blast-fragmentation types depending on the target set. They are the standard warhead type for the vast majority of tactical missiles and guided munitions in service worldwide. Warhead design must be matched to the target vulnerability, fuze mode, and engagement geometry to maximize the probability of kill against the intended target class.}{catRM}

\dictentry{Countdown}{The timeline leading to launch, a synchronized sequence of events and system checks that proceeds from propellant loading through engine ignition to liftoff. The countdown provides structured milestones for propellant loading, tank pressurization, guidance alignment, and engine health verification before committing to launch. A hold or recycle capability allows the countdown to be paused or restarted if an anomaly is detected before T-0. The countdown discipline ensures that every critical system is verified in the correct sequence and that the vehicle is ready for safe flight.}{catRM}

\dictentry{Countermeasure}{A system defeating enemy targeting by disrupting, deceiving, or degrading the threat sensor or guidance system. Countermeasures include chaff, flares, electronic jammers, decoys, and low-observable technologies that reduce the probability of detection or successful engagement. Effective countermeasures are continuously evolving in response to improvements in threat sensor and guidance technology. The selection and deployment of countermeasures depends on the threat environment, platform capabilities, and the specific sensor types that the adversary employs for targeting.}{catRM}

\dictentry{Cruise Missile}{A self-navigating missile with continuous thrust that flies at sustained subsonic or supersonic speed toward its target using inertial, GPS, or terrain-matching guidance. Unlike ballistic missiles, cruise missiles maintain powered flight throughout their trajectory and can maneuver to avoid defenses and follow terrain-hugging profiles. Examples include the Tomahawk and Kalibr, which are launched from ships, submarines, and aircraft for precision strike missions. Their low-altitude flight profiles make them difficult to detect by ground-based radar and air defense systems.}{catRM}

\dictentry{Cryogenic Propellant}{Propellant stored at extremely low temperatures, most commonly liquid hydrogen at -253 degrees C or liquid oxygen at -183 degrees C. Cryogenic propellants offer high specific impulse due to their low molecular weight and high combustion energy, but require insulated tanks and continuous boil-off management. They are used on high-performance launch vehicles such as the Space Shuttle main engines, Delta IV, and SLS where maximum performance outweighs the operational complexity of cryogenic handling. Ground support equipment for cryogenic propellants includes vacuum-jacketed transfer lines, conditioning systems, and real-time tanking instrumentation.}{catRM}

\dictentry{David Sling}{An Israeli medium-range air defense system also known as Magic Wand, designed to intercept tactical ballistic missiles, cruise missiles, and large-caliber rockets. It uses the Stunner interceptor with a unique kinetic kill warhead that combines electro-optical and radar seekers for precise target discrimination. David Sling fills the gap between the Iron Dome short-range system and the Arrow upper-tier ballistic missile defense system. The system entered operational service in 2017 and provides Israel with a layered defense capability against a wide spectrum of aerial threats.}{catRM}

\dictentry{De Laval Nozzle}{A convergent-divergent nozzle for supersonic exhaust that accelerates subsonic combustion gases through a sonic throat and then expands them supersonically in a divergent section. The nozzle converts the thermal energy of the combustion products into directed kinetic energy, producing thrust. The de Laval nozzle design is fundamental to all high-performance rocket engines and was first developed in the 19th century by Gustaf de Laval. The throat-to-exit area ratio determines the final exhaust velocity and must be optimized for the intended operating altitude.}{catRM}

\dictentry{Decoy}{A device misleading enemy defenses by presenting a false radar, infrared, or visual signature that mimics the real threat. Decoys can be passive reflectors, active repeaters, or expendable flares that draw defensive fire and sensors away from the actual warhead or aircraft. In ballistic missile defense, decoys complicate the discrimination problem by forcing interceptors to distinguish real warheads from numerous decoys during midcourse. The effectiveness of decoys depends on their ability to closely replicate the signature characteristics of the real threat across multiple sensor bands.}{catRM}

\dictentry{Deorbit Burn}{A burn to initiate return to Earth performed in the direction opposite to the orbital velocity vector to reduce orbital energy. The retrograde burn lowers the perigee into the atmosphere, allowing aerodynamic drag to decelerate the vehicle for reentry at the desired location. Deorbit burns are precisely timed and computed to achieve the target landing site within the recovery or impact zone. The magnitude of the deorbit burn determines the reentry angle and directly affects the thermal and mechanical loads experienced during atmospheric entry.}{catRM}

\dictentry{Depressed Trajectory}{A flatter than normal ballistic trajectory that reduces the flight time and peak altitude of a ballistic missile. Depressed trajectories sacrifice some range to arrive at the target faster, reducing the time available for enemy detection and interception. This technique is particularly valuable for penetrating missile defense systems by minimizing the midcourse phase when the warhead is most vulnerable to intercept. The tradeoff between reduced flight time and decreased range must be carefully balanced against mission requirements and threat environment.}{catRM}

\dictentry{Destruct System}{A system destroying the rocket in flight for range safety, consisting of explosive charges or ordnance that sever the motor case and terminate thrust if the vehicle deviates from its planned trajectory. Ground-based range safety officers monitor the flight path and command destruction if the missile poses a danger to populated areas. The destruct system is an essential safety requirement for all launch vehicles operating over public lands or waters. Redundant command receivers and autonomous flight termination logic ensure that the system can function even if ground communication is lost.}{catRM}

\dictentry{Discharge Line}{A line on the outlet side of the pump that carries high-pressure propellant from the turbopump discharge to the main propellant valve and ultimately to the injector. The discharge line must withstand full operating pressure and accommodate thermal expansion, vibration, and misalignment between the pump and engine. Its diameter and routing are designed to minimize pressure losses and avoid hydraulic resonance that could contribute to system instability. Proper discharge line design ensures smooth, pulse-free propellant delivery to the injector under all operating conditions.}{catRM}

\dictentry{Drain Valve}{A valve for draining propellant from tanks and lines during ground operations, maintenance, and safing procedures. Drain valves must provide leak-tight closure under operating pressures and be compatible with the propellant chemicals without causing contamination or corrosion. They are typically located at the lowest points of the propellant system to ensure complete drainage by gravity. Drain valve design must also account for safe handling of toxic or cryogenic propellants during ground servicing operations.}{catRM}

\dictentry{Droplet Combustion}{Burning of individual propellant droplets as they travel through the combustion chamber, where the flame front surrounds each droplet and vaporized fuel reacts with oxidizer. The droplet combustion rate is governed by the D-squared law, where the square of droplet diameter decreases linearly with time. Understanding droplet combustion is essential for optimizing injector design and predicting combustion efficiency in liquid rocket engines. The initial droplet size distribution produced by the injector directly determines the combustion completeness and overall engine performance.}{catRM}

\dictentry{Dry Mass}{The mass of the rocket without propellant, including structure, engines, avionics, and payload. Dry mass is a critical parameter in the rocket equation because minimizing it directly increases the delta-v available for the mission. Structural mass fraction improvements through advanced materials and design techniques are a primary focus of launch vehicle development programs. Every kilogram saved in dry mass translates directly into additional payload capacity or extended mission range for the launch vehicle.}{catRM}

\dictentry{Dual Bell Nozzle}{A nozzle with two bell contours for altitude compensation that operates as a conventional nozzle at low altitude and transitions to an extended expansion mode at high altitude. The inflection point in the nozzle contour allows ambient pressure to control the flow separation, automatically adapting the expansion ratio. This passive altitude compensation concept offers a simpler alternative to mechanically deployable nozzle extensions. The dual bell design improves overall mission performance by optimizing expansion efficiency across both sea-level and vacuum operating conditions.}{catRM}

\dictentry{Dual Thrust Motor}{A motor with two thrust levels, typically a high-thrust boost phase followed by a lower-thrust sustain phase within a single motor case. The grain geometry is designed with different burning surface areas to produce the required thrust profile without complex valve or control mechanisms. Dual thrust motors are widely used in air-to-air missiles and tactical rockets where rapid acceleration followed by extended range is needed. The transition between boost and sustain phases must be smooth to avoid structural loads or guidance disruption during flight.}{catRM}

\dictentry{Early Warning Radar}{Radar providing early threat detection by scanning large volumes of airspace for incoming ballistic missiles, aircraft, or cruise missiles. Early warning radars such as the AN/FPS-132 and PAVE PAWS use powerful transmitters and large apertures to detect targets at ranges of thousands of kilometers. Their data feeds missile defense interceptors and national command authorities to enable timely defensive response. The radar must reliably detect and track targets while rejecting clutter, decoys, and electronic countermeasures in a complex threat environment.}{catRM}

\dictentry{Effective Exhaust Velocity}{The effective speed of exhaust including pressure effects, calculated as specific impulse multiplied by standard gravity. It represents the equivalent velocity that, if all exhaust were expelled at this speed, would produce the same thrust as the actual non-uniform exhaust. Effective exhaust velocity is used in the rocket equation to compute delta-v and is directly related to engine performance. Higher effective exhaust velocity means more efficient propellant utilization and greater payload delivery capability for a given vehicle mass.}{catRM}

\dictentry{Electric Pump Cycle}{A cycle using electric motors to drive propellant pumps instead of gas turbines, eliminating the need for a gas generator or preburner. The electric pumps are powered by batteries or other electrical sources, offering independent pump speed control and simplifying the engine cycle. This concept is used on engines such as the Rocket Lab Rutherford, where the reduced complexity and part count benefit small launch vehicle economics. The elimination of hot gas turbomachinery reduces development risk and enables faster engine prototyping and iteration.}{catRM}

\dictentry{Electronic Countermeasure}{Electronic jamming or deception used to disrupt enemy radar, communications, and missile guidance systems. ECM techniques include noise jamming to saturate receivers, deceptive jamming to create false targets, and cross-eye techniques to distort angle tracking. Modern ECM systems are adaptive, sensing the threat environment in real time and selecting optimal countermeasure techniques to protect the platform. Effective ECM requires detailed knowledge of the threat radar characteristics and continuous updating as adversaries develop new waveforms and counter-countermeasure techniques.}{catRM}

\dictentry{End Burning}{Burning from one end of the grain in a longitudinal direction, producing a steadily decreasing burning surface area as the flame front recedes. This grain configuration yields a regressive thrust profile and is used when a long, low-thrust burn is required such as in sustainer motors. End burning grains are simple to manufacture but offer less design flexibility for thrust tailoring compared to more complex grain geometries. The burn rate is directly controlled by propellant formulation, providing predictable and repeatable performance across production lots.}{catRM}

\dictentry{Engine Cycle}{The propellant feed cycle of a liquid rocket engine, defining how propellant is pressurized and delivered to the combustion chamber. Common cycles include gas generator, staged combustion, expander, tap-off, and pressure-fed, each offering different tradeoffs between performance, complexity, and reliability. The choice of engine cycle is one of the most consequential design decisions in liquid rocket engine development. The cycle determines the achievable chamber pressure, specific impulse, engine weight, and the overall complexity of the propulsion system.}{catRM}

\dictentry{Erosion}{Material loss from the nozzle throat during firing caused by high-temperature gas flow, chemical reactions, and particulate impact on the throat insert surface. Throat erosion increases the throat area over time, which reduces chamber pressure and degrades performance. Throat inserts are made from refractory materials such as graphite, carbon-carbon, or tungsten to minimize erosion rate during the burn. Predicting and controlling erosion is critical for long-duration burns where accumulated throat area growth can significantly affect the thrust profile and mission performance.}{catRM}

\dictentry{Exhaust Plume}{The visible exhaust of a rocket engine consisting of hot combustion gases, particulates, and condensates expanding into the ambient atmosphere. The plume shape and brightness depend on chamber pressure, ambient pressure, propellant type, and expansion ratio. Plume analysis is important for predicting infrared signatures, assessing thermal loads on the vehicle base, and modeling exhaust contamination on the launch pad. Plume characteristics also affect the electromagnetic signature of the vehicle and can influence radar detection and tracking during boost phase.}{catRM}

\dictentry{Exhaust Trail}{The visible trail from a missile exhaust composed of condensed water vapor, soot, and ice crystals that persist in the atmosphere behind the vehicle. Exhaust trails can reveal a missile's position and trajectory to visual observers and tracking systems long after the missile has passed. In some cases, exhaust trails are studied meteorologically to assess the impact of rocket launches on upper-atmosphere chemistry. The persistence and visibility of exhaust trails depend on atmospheric conditions, altitude, and the propellant combustion products.}{catRM}

\dictentry{Exhaust Velocity}{The speed of exhaust gases leaving the nozzle relative to the engine, determined by the ratio of combustion temperature to exhaust molecular weight and the nozzle expansion ratio. Higher exhaust velocity directly translates to greater specific impulse and more efficient use of propellant. Maximizing exhaust velocity while managing heat loads and material constraints is a central challenge in rocket engine design. The exhaust velocity varies across the nozzle exit plane, with the highest velocities occurring along the centerline and decreasing toward the wall boundary layer.}{catRM}

\dictentry{Expander Cycle}{A cycle using heat absorption to drive the turbine, where one of the propellants is heated by the combustion chamber or nozzle walls before being routed through the turbine to drive the propellant pump. The expander cycle provides a simple and reliable engine architecture with no gas generator and is well-suited to upper-stage applications. Notable examples include the RL10 and Vinci engines, which have accumulated extensive flight heritage. The limited heat transfer available restricts expander cycle engines to moderate thrust levels, making them ideal for upper stages rather than booster applications.}{catRM}

\dictentry{Expansion Ratio}{The ratio of nozzle exit area to throat area, which determines how much the exhaust gases are expanded and thus how efficiently thermal energy is converted to thrust. A higher expansion ratio produces more thrust in a vacuum but can cause flow separation at sea level where ambient pressure is higher. Optimal expansion ratio depends on the operating altitude and is a key design parameter for nozzle sizing. Single-expansion-ratio nozzles represent a compromise between sea-level and vacuum performance that must be optimized for the intended mission profile.}{catRM}

\dictentry{Expansion Ratio Effect}{The effect of nozzle expansion on performance, where underexpanded nozzles lose thrust because exhaust exits at higher pressure than ambient, while overexpanded nozzles risk flow separation and lateral loads on the nozzle wall. The ideal design point occurs at optimal expansion where exit pressure matches ambient pressure. Multi-stage and altitude-compensating nozzles attempt to mitigate the expansion ratio effect across the flight envelope. Understanding this effect is essential for nozzle designers seeking to maximize performance over the full range of operating altitudes encountered during a mission.}{catRM}

\dictentry{Expendable Launch Vehicle}{A single-use launch vehicle that is discarded after delivering its payload to orbit, with no recovery or reuse of any stages. ELVs are simpler and often lighter than reusable designs because they do not require recovery hardware such as heat shields, landing gear, or restart-capable engines. Despite higher per-launch costs, ELVs remain widely used for their reliability and operational simplicity. The absence of reuse constraints allows ELV structures to be optimized purely for ascent loads, potentially reducing dry mass compared to reusable counterparts.}{catRM}

\dictentry{Fairing Separation}{Jettisoning the payload fairing after atmospheric exit once aerodynamic heating and pressure loads become negligible. The fairing is split into halves or petals and pushed away by pyrotechnic or pneumatic mechanisms to expose the payload for deployment. Fairing separation must occur as early as possible to reduce dead weight carried to orbit while ensuring the payload is not exposed to damaging dynamic pressures. The separation event must produce clean, symmetric deployment to avoid recontact with the payload or upper stage.}{catRM}

\dictentry{Fiber Optic Guidance}{Guidance via fiber optic link where an unspooling fiber connects the missile to the launch platform, transmitting real-time targeting data and imagery. The operator can observe the target area through a seeker camera and send steering corrections over the fiber, enabling precision engagement of moving or partially obscured targets. Fiber optic guidance is used in systems like the EFOG-M and Spike NLOS, offering high resistance to electronic jamming. The fiber link provides a secure, high-bandwidth data connection that cannot be intercepted or jammed by electromagnetic means.}{catRM}

\dictentry{Fill Valve}{A valve for filling propellant tanks during ground operations, designed to allow controlled flow of propellant into the tank while preventing overfill and managing venting of displaced gas. Fill valves are sized for rapid propellant loading within countdown timeline constraints and must provide leak-tight closure during flight. They are part of the ground support equipment interface and are typically closed and safed before vehicle lift-off. Fill valve design must account for thermal contraction when cryogenic propellants are loaded and must prevent vapor lock during high-flow filling operations.}{catRM}

\dictentry{Film Cooling}{Cooling by injecting a coolant film on the chamber wall, where a thin layer of cooler propellant or inert fluid flows along the inner surface to insulate it from hot combustion gases. Film cooling reduces thermal loads on the chamber wall but sacrifices some performance because the coolant film does not fully participate in combustion. The balance between wall protection and performance loss is a key tradeoff in regeneratively cooled engine design. The number of film cooling injector rows and the coolant flow fraction are optimized during development to achieve adequate protection with minimal performance penalty.}{catRM}

\dictentry{Filter}{A device removing particles from propellant to prevent clogging of injector orifices, damage to pump seals and bearings, and contamination of the combustion chamber. Filters are rated by their mesh size or micrometer rating and are placed at key locations throughout the propellant feed system. Regular inspection and replacement of filters is an important maintenance practice to ensure reliable engine operation. Filter integrity is verified before each flight through pressure drop measurements and visual inspection to confirm that no bypass or structural failure has occurred.}{catRM}

\dictentry{Fin-Stabilized Rocket}{A rocket stabilized by aerodynamic fins mounted on the rear that generate restoring forces and moments when the rocket departs from its velocity vector. The fins shift the center of pressure behind the center of gravity, creating a natural stability that keeps the nose pointed into the relative wind. Fin-stabilized rockets are simple and inexpensive, making them common in unguided artillery rockets and small sounding rockets. The fin size, shape, and sweep angle are designed to provide adequate stability margin across the full range of flight velocities and angles of attack.}{catRM}

\dictentry{Fire Control Radar}{Radar directing weapons by providing precise range, bearing, and velocity data for targeting calculations and missile guidance. Fire control radars typically operate in higher frequency bands with narrow beam widths for accurate target tracking and illumination. They are integrated with fire control computers that generate firing solutions and guide semi-active or command-guided missiles to their targets. The radar must maintain continuous target track quality while operating in cluttered and contested electromagnetic environments with active enemy countermeasures.}{catRM}

\dictentry{First Stage}{The initial stage providing lift-off thrust, which is the largest and most powerful section of a multi-stage rocket responsible for lifting the entire vehicle off the launch pad and through the dense lower atmosphere. The first stage typically burns for 2 to 3 minutes and delivers the vehicle to an altitude of 50 to 80 kilometers before separation. It uses the most powerful engines on the rocket and may include strap-on boosters for additional thrust. First stage design must balance high thrust for liftoff against the structural mass required to support the full vehicle weight during ground operations and initial ascent.}{catRM}

\dictentry{Flame Deflector}{A structure deflecting exhaust at launch to redirect hot gases away from the launch pad, vehicle, and ground support equipment. Flame deflectors are typically water-cooled steel or concrete structures positioned beneath the launch mount. They prevent reflected exhaust from damaging the vehicle and reduce acoustic loading during the critical first seconds of flight. The deflector geometry must handle extreme thermal loads and high-velocity exhaust impingement while channeling gases safely into the flame trench below the pad.}{catRM}

\dictentry{Flare}{A heat source decoying infrared seekers by presenting a higher-intensity thermal signature than the aircraft or missile it is protecting. Flares are ejected from dispensers in programmed sequences to create a cluster of intense heat sources that seduce IR-guided missiles away from the target. Modern advanced IR seekers use multi-color detection and tracking algorithms to discriminate between flares and real targets, driving the development of more sophisticated flare countermeasures. Advanced flare compositions and ejection patterns continue to evolve to maintain effectiveness against increasingly capable infrared missile threats.}{catRM}

\dictentry{Flexible Bearing}{A flexible joint for nozzle gimballing that allows the thrust vector to be steered by permitting angular deflection of the nozzle while maintaining a pressure-tight seal. Flexible bearings consist of an elastomeric laminate structure that flexes under actuator loads while containing combustion pressure. They are lighter and simpler than ball-and-socket joints and are used on engines like the SSME and Merlin for thrust vector control. The bearing must withstand millions of pressure cycles and thermal extremes while maintaining a reliable seal throughout the engine operational life.}{catRM}

\dictentry{Flight Termination System}{A system for range safety destruct that uses redundant ordnance charges and command receivers to destroy the vehicle if it deviates from the planned trajectory. The FTS is independently powered, uses separation and safety sensors, and can be triggered by ground command or autonomously by onboard logic. It is a mandatory safety system for all military and commercial launch vehicles operating in controlled airspace. The FTS must function reliably even after sustaining damage from vehicle breakup or other catastrophic events during flight.}{catRM}

\dictentry{Flow Separation in Nozzle}{Detachment of flow from the nozzle wall occurring when the ambient pressure exceeds the local gas pressure at the nozzle wall, typically in overexpanded nozzles at low altitude. Flow separation causes asymmetric lateral loads, reduced performance, and potential structural damage to the nozzle. Nozzle designers account for separation through contour optimization, boundary layer management, and operational altitude constraints. Predicting the exact onset of flow separation requires detailed computational fluid dynamics analysis combined with experimental validation during engine testing.}{catRM}

\dictentry{Fragmentation Warhead}{A warhead releasing fragments that generates a cloud of high-velocity metal shards upon detonation to damage the target through kinetic energy impact. The casing is pre-scored or incorporates notched liners to produce controlled fragment sizes optimized for the target type. Fragmentation warheads are effective against aircraft, light vehicles, and personnel, and are the most common warhead type in surface-to-air and air-to-air missiles. The fragment velocity, mass distribution, and spatial pattern are carefully engineered to maximize the probability of kill within the lethal radius of the warhead.}{catRM}

\dictentry{Full Flow Staged Combustion}{A staged combustion cycle with two preburners that gasify all of both propellants before they enter the main combustion chamber, with one preburner running fuel-rich and the other oxidizer-rich. Every drop of propellant passes through the turbine, maximizing the energy extracted and enabling very high chamber pressures and specific impulse. This demanding cycle is used on the SpaceX Raptor engine, achieving performance levels not possible with gas generator or conventional staged combustion cycles. The full flow architecture eliminates the turbine blade cooling problem since all gas entering the turbine is above the melting point of the blade material.}{catRM}

\dictentry{Gas Cooler}{A device cooling gas before use, typically to protect turbine blades and bearings from the extreme temperatures of combustion gases. Heat exchangers, quench chambers, or secondary fluid injection can reduce gas temperature to within material limits while preserving enough energy to drive the turbopump. Gas cooler effectiveness directly impacts turbine life and the overall reliability of gas-generator-cycle engines. The cooler must achieve sufficient temperature reduction without introducing excessive pressure drop that would reduce the energy available for turbine work.}{catRM}

\dictentry{Gas Generator}{A device generating gas to drive a turbine by burning a small portion of the main propellants in a fuel-rich or oxidizer-rich combustion process. The resulting warm gas expands through the turbine to spin the propellant pumps before being dumped overboard or routed to the nozzle. Gas generators are compact and relatively simple, forming the heart of gas-generator-cycle engines like the F-1 and Merlin. The gas generator must produce stable, clean combustion across a wide range of operating conditions to ensure reliable turbopump operation throughout the engine burn.}{catRM}

\dictentry{Gas Generator Cycle}{A cycle burning a small portion of propellant to drive a turbine, with the turbine exhaust dumped overboard or into the nozzle at suboptimal pressure. This cycle offers simplicity and moderate performance at the cost of some wasted propellant energy. It is the most common engine cycle for first-stage launch vehicle engines, providing a good balance of thrust, reliability, and development cost. The gas generator cycle is favored for booster applications where the simplicity advantage outweighs the performance penalty relative to more complex staged combustion cycles.}{catRM}

\dictentry{Gimbal}{A swivel mount for engine thrust vectoring that allows the engine to pivot and redirect the thrust vector relative to the vehicle centerline. Gimballing is controlled by hydraulic or electromechanical actuators receiving commands from the flight control system. Thrust vector control through gimbaling provides the primary steering mechanism for liquid rocket engines during powered flight. The gimbal bearing must support the full thrust load while permitting smooth angular deflection under the combined effects of acceleration, vibration, and thermal growth.}{catRM}

\dictentry{GPS Guidance}{Guidance using satellite positioning from the Global Positioning System constellation to determine the missile's position, velocity, and time in real time. GPS guidance provides high accuracy that is independent of range and does not degrade with distance, making it ideal for precision strike munitions and cruise missiles. GPS-aided inertial navigation systems combine satellite fixes with inertial sensors to produce robust and accurate guidance throughout the flight. GPS guidance can be degraded by intentional jamming or spoofing, which is why it is typically combined with inertial navigation as a backup.}{catRM}

\dictentry{Grain}{A shaped solid propellant charge molded or cast into a specific geometry within the motor case to control the burning surface evolution over time. The grain shape determines the thrust-time profile of the motor, with different geometries producing progressive, neutral, or regressive thrust curves. Grain design is a critical element of solid rocket motor engineering, directly affecting performance, safety, and structural integrity. The grain must maintain its structural integrity under thermal cycling, transportation loads, and the inertial forces experienced during motor operation without cracking or debonding.}{catRM}

\dictentry{Grain Configuration}{The geometric shape of the solid grain, including its cross-sectional profile and axial features such as slots, perforations, and inhibitors. Common configurations include circular bore, star, finocyl, and wagon wheel, each producing a distinct thrust-time history. Grain configuration is the primary design variable for tailoring the performance of solid rocket motors to meet mission-specific thrust requirements. The selected configuration must be compatible with the casting and curing manufacturing process while providing the required burning surface area throughout the burn duration.}{catRM}

\dictentry{Graphite Nozzle}{A nozzle using graphite for high temperature resistance, capable of withstanding combustion temperatures above 3000 K with low erosion rates. Graphite is machined to precise throat and contour dimensions and is often used as a throat insert in solid rocket motors and short-duration liquid engines. Isostatically pressed graphite grades offer the best combination of thermal shock resistance and mechanical strength for rocket nozzle applications. Graphite nozzles are cost-effective for single-use applications but are less suitable for reusable engines due to oxidation and erosion over multiple burns.}{catRM}

\dictentry{Ground-Based Interceptor}{A US ground-based missile interceptor deployed at Fort Greely in Alaska and Vandenberg Air Force Base in California as the primary weapon of the Ground-based Midcourse Defense system. The GBI uses a kinetic kill vehicle that destroys incoming ballistic missile warheads by direct body-to-body impact in space. GBIs are the only US interceptors currently capable of engaging intermediate-range and intercontinental ballistic missiles during the midcourse phase. Each GBI carries a single exoatmospheric kill vehicle and is housed in a silo-based launch canister ready for rapid launch on warning.}{catRM}

\dictentry{Guidance System}{The system directing a missile to its target, integrating sensors, computers, and control actuators to navigate the missile along a computed trajectory. Guidance systems can be inertial, GPS-aided, radar-homing, infrared-homing, laser-homing, or command-guided depending on the missile type and engagement scenario. The accuracy and responsiveness of the guidance system is the primary determinant of a missile's single-shot probability of kill. Modern guidance systems integrate multiple sensor modalities and adaptive algorithms to maintain accuracy in the presence of countermeasures and environmental disturbances.}{catRM}

\dictentry{Heat Exchanger}{A device transferring heat between fluids, used in rocket engines to warm pressurant gas, preheat propellant, or cool hot components. In expander-cycle engines, the heat exchanger absorbs combustion heat through the chamber wall to generate turbine drive gas. Heat exchanger design must balance thermal transfer efficiency with pressure drop, weight, and structural integrity under high thermal gradients. The effectiveness of the heat exchanger directly impacts engine cycle performance, turbine inlet temperature, and overall system efficiency.}{catRM}

\dictentry{Helium Pressurization}{Pressurization using helium, an inert gas with low molecular weight that does not react with propellants and remains gaseous at cryogenic temperatures. Helium is stored in high-pressure composite overwrapped tanks and regulated to the required tank pressure before being introduced above the propellant. This is the most common pressurization method for liquid rocket engines, providing reliable and contamination-free tank pressurization. The helium supply must be sized to maintain tank pressure throughout the burn as propellant volume decreases and thermal conditions change.}{catRM}

\dictentry{Hit-to-Kill}{Direct impact interception where the interceptor destroys the target by kinetic energy impact at combined closing velocities of several kilometers per second. No explosive warhead is needed because the sheer kinetic energy of the collision shatters both vehicles. Hit-to-kill requires extreme guidance precision, as the interceptor must achieve meter-level accuracy at intercept, and is used in systems like THAAD, GBI, and SM-3. The kinetic energy at typical intercept velocities is equivalent to several tons of TNT, ensuring complete destruction of the target warhead.}{catRM}

\dictentry{Hot Gas Valve}{A valve directing hot gas for thrust vectoring or gas management, operating at extreme temperatures and pressures without external cooling. Hot gas valves use heat-resistant materials and specialized sealing designs to maintain function in the harsh environment of rocket exhaust or turbine discharge. They are critical components in gas-generator and tap-off cycle engines where hot gas must be precisely controlled for turbine drive or dump. These valves must operate reliably through thousands of thermal cycles while maintaining leak-tight closure under extreme pressure differentials.}{catRM}

\dictentry{Hybrid Rocket}{A rocket combining solid fuel with liquid or gaseous oxidizer that offers simplicity and safety advantages over both solid and liquid rockets. The solid fuel grain cannot detonate or explode on its own, and the oxidizer flow can be throttled or shut off to control or terminate thrust. Hybrid rockets are used in sounding rockets, launch vehicle upper stages, and emerging commercial spaceflight applications where safety and low cost are priorities. The regression rate of the solid fuel surface limits thrust levels, making hybrids better suited to moderate-thrust applications than high-thrust booster stages.}{catRM}

\dictentry{Hydrazine}{A storable monopropellant and fuel with the chemical formula N2H4 that decomposes exothermically over an iridium catalyst to produce hot gas for thrust. Hydrazine is also used as a fuel in bipropellant combinations with nitrogen tetroxide, offering hypergolic ignition and long-term storability. Despite its toxicity and handling hazards, hydrazine remains widely used in spacecraft attitude control, orbital maneuvering, and missile propulsion systems. Safe handling requires specialized protective equipment, closed-loop transfer systems, and dedicated storage facilities to protect personnel and the environment from exposure.}{catRM}

\dictentry{Hypergolic}{A propellant combination that ignites on contact without an external ignition source, providing highly reliable and instant combustion when the two components are mixed. Common hypergolic pairs include monomethylhydrazine with nitrogen tetroxide and UDMH with nitrogen tetroxide. Hypergolic propellants are storable at ambient conditions, making them ideal for missile and spacecraft propulsion where long shelf life and rapid launch capability are essential. The spontaneous ignition reliability of hypergolic propellants eliminates the need for complex ignition systems and significantly simplifies engine design and ground operations.}{catRM}

\dictentry{Hypersonic Glide Vehicle}{A maneuverable hypersonic boost-glide weapon that is first boosted to high altitude on a ballistic or near-ballistic trajectory, then glides through the upper atmosphere at speeds exceeding Mach 5 on a lifting trajectory. HGVs can perform cross-range and down-range maneuvers that make their impact point unpredictable and extremely difficult to intercept. Systems like the Chinese DF-ZF and Russian Avangard represent a new class of strategic weapons that challenge existing missile defense architectures. The extreme thermal environment during hypersonic glide requires advanced thermal protection materials and aerodynamic shaping to maintain structural integrity throughout the flight.}{catRM}

\dictentry{Hypersonic Missile}{A missile capable of speeds above Mach 5 that maintains sustained powered or glide flight through the atmosphere at extreme velocity. Hypersonic missiles combine the speed of ballistic weapons with the maneuverability of cruise missiles, drastically reducing reaction time for defenders. They may use scramjet, ramjet, or boost-glide propulsion, and represent a rapidly developing class of weapons among major military powers. The extreme aerodynamic heating at hypersonic speeds demands advanced thermal protection systems and materials that can withstand sustained exposure to temperatures exceeding several thousand degrees.}{catRM}

\dictentry{ICBM}{An intercontinental ballistic missile with a range exceeding 5500 kilometers capable of delivering nuclear or conventional warheads between continents. ICBMs follow a ballistic trajectory through space with a midcourse phase lasting approximately 20 to 30 minutes, during which the reentry vehicle coasts outside the atmosphere. They are a primary strategic deterrent carried by land-based silos, mobile launchers, and submarine-launched variants. Modern ICBMs carry multiple independently targetable reentry vehicles and penetration aids to ensure survivability against advanced missile defense systems.}{catRM}

\dictentry{Igniter}{A device initiating combustion in the chamber by providing an initial heat source sufficient to raise the propellant to its autoignition temperature. Igniters can be pyrotechnic, spark-based, catalytic, or pyrogen, with the choice depending on propellant type and engine requirements. Reliable ignition is critical for mission success, and igniter systems are typically designed with redundancy to ensure startup even if one component fails. The igniter must deliver sufficient energy to establish stable combustion across the entire injector face within the required start transient timeframe.}{catRM}

\dictentry{Igniter Grain}{A small grain for motor ignition that burns rapidly and produces hot gas to ignite the main propellant charge. The igniter grain is typically a pyrotechnic composition pressed into a small cylindrical or tubular shape and activated by an electrical squib. It provides the initial energy pulse that establishes the combustion wave across the main grain surface for uniform and reliable motor startup. The igniter grain formulation and geometry are carefully designed to produce the correct pressure-time trace for safe and predictable motor ignition.}{catRM}

\dictentry{Impact Point}{The point where the missile strikes the target or ground, determined by the trajectory, guidance accuracy, and terminal maneuvering capability. The predicted impact point is continuously computed by the guidance system and compared to the aim point to generate steering corrections. Impact point prediction is also used by missile defense systems to determine the predicted intercept point for deploying interceptors. Accurate impact point prediction requires real-time knowledge of the missile's state vector, atmospheric conditions, and the gravitational field along the trajectory.}{catRM}

\dictentry{Impinging Injector}{An injector where jets impinge for mixing by directing opposing streams of fuel and oxidizer to collide at a designated point. The collision atomizes and mixes the propellants into a fine spray before they enter the combustion zone, promoting rapid and complete combustion. Impinging injectors are used in many pressure-fed and gas-generator-cycle engines where simplicity and good atomization are required. The impingement point location and the angle between streams are critical design parameters that determine mixing quality and combustion stability.}{catRM}

\dictentry{Inertial Guidance}{Guidance using accelerometers and gyroscopes to continuously measure the missile's acceleration and orientation, computing position and velocity through dead-reckoning integration. Inertial guidance is self-contained and requires no external references, making it immune to jamming and applicable to any flight environment. Accuracy degrades over time due to sensor drift, so inertial systems are often supplemented with GPS or stellar updates for long-range missiles. The quality of the inertial sensors, particularly gyro bias stability and accelerometer linearity, directly determines the navigation accuracy achievable over the mission duration.}{catRM}

\dictentry{Infrared Homing}{Homing using heat emissions from the target, where an onboard infrared seeker detects the thermal radiation of the target aircraft's engine exhaust or aerodynamic heating and steers toward the strongest source. IR homing provides passive, fire-and-forget engagement capability with no warning to the target. Modern imaging infrared seekers can distinguish target features and are resistant to flares and other heat-based countermeasures. Advanced IR seekers use focal plane arrays with thousands of detector elements to achieve high angular resolution and target discrimination capability.}{catRM}

\dictentry{Inhibitor}{A coating preventing burning on certain grain surfaces by using a material that resists combustion and insulates that surface from the flame front. Inhibitors are applied to specific faces of the grain to control which surfaces burn, and thus define the burning surface evolution and resulting thrust profile. Common inhibitor materials include rubber compounds, epoxy resins, and phenolic coatings that char but do not propagate the combustion front. The bond between the inhibitor and the propellant grain must remain intact throughout thermal cycling and motor operation to prevent unintended surface exposure.}{catRM}

\dictentry{Injection Burn}{A burn to inject into a transfer trajectory, typically a Hohmann transfer or bi-elliptic maneuver, to move from one orbit to another. The injection burn changes the velocity of the spacecraft to place it on an elliptical path whose apogee or perigee reaches the target orbit. Precise timing and delta-v magnitude are critical for achieving the correct orbit without requiring excessive correction maneuvers. The injection burn is typically the most energy-demanding maneuver in an orbital mission and determines the overall success of the orbit transfer.}{catRM}

\dictentry{Injector}{A device mixing fuel and oxidizer in the combustion chamber by spraying propellants through precision orifices in a carefully designed pattern. Injector design controls atomization quality, spray distribution, propellant impingement patterns, and mixing efficiency, all of which directly affect combustion performance and stability. The injector face is one of the most thermally stressed components in a liquid rocket engine. Injector development involves extensive cold-flow and hot-fire testing to validate spray characteristics, combustion stability, and thermal loads before flight qualification.}{catRM}

\dictentry{Injector Pattern}{The arrangement of injector elements across the face plate that determines the spatial distribution of fuel and oxidizer streams entering the chamber. Common patterns include concentric ring, showerhead, and self-impinging designs, each offering different tradeoffs in mixing, stability, and manufacturing complexity. The injector pattern is a primary design variable for controlling combustion stability and heat release distribution. The pattern must provide uniform propellant distribution while avoiding coherent acoustic modes that could couple with combustion and drive instability.}{catRM}

\dictentry{Insertion}{The process of entering the intended orbit by executing one or more propulsive burns at precise points along the trajectory. Orbit insertion requires matching the target orbit's altitude, inclination, and eccentricity within tight tolerances. The insertion burn is typically performed by the upper stage or spacecraft propulsion system after the launch vehicle has delivered the payload to the appropriate transfer orbit. Multiple burns may be required for complex orbit geometries, with each burn precisely timed and sized to achieve the desired orbital parameters.}{catRM}

\dictentry{Insulation}{Thermal protection inside the motor case that prevents combustion heat from reaching the case wall and causing structural failure. Insulation layers are applied to the case bore, head end, and nozzle interfaces using elastomeric or phenolic materials that char and ablate during operation. Proper insulation design is critical for motor safety, as premature case burnthrough can result in catastrophic failure. The insulation thickness is tailored to the expected heat flux profile, with thicker layers in regions of high thermal exposure such as near the nozzle and head end.}{catRM}

\dictentry{Interceptor}{A missile designed to destroy incoming threats by direct impact or proximity-fuzed warhead detonation. Interceptors are optimized for high speed, rapid response, and precise terminal guidance to engage ballistic missiles, cruise missiles, aircraft, or rockets. They are the kinetic element of missile defense systems and are deployed in ground-based, sea-based, and mobile configurations. Interceptor design requires balancing speed, maneuverability, and seeker performance against the constraints of weight, size, and reaction time from detection to launch.}{catRM}

\dictentry{Internal Burning}{Burning from internal passages within the grain, where the combustion front propagates outward from perforations or slots in the grain interior. Internal burning allows the motor case to be structurally loaded by chamber pressure, making the case a true pressure vessel. This grain configuration is standard for most solid rocket motors because it provides design flexibility and predictable structural behavior. The internal burning surface evolves in a controlled manner determined by the perforation geometry, enabling precise thrust profile tailoring.}{catRM}

\dictentry{Interstage}{A structure connecting rocket stages that provides structural continuity during powered flight and houses the separation mechanism, avionics, and wiring between stages. The interstage must be strong enough to transmit axial thrust loads and bending moments while being light enough to minimize the performance penalty of carried dead weight. Interstage structures are typically jettisoned with the spent lower stage to reduce the mass carried to orbit. The interstage design must accommodate thermal environments from both the operating engine below and the vacuum conditions above.}{catRM}

\dictentry{IRBM}{An intermediate-range ballistic missile with a range between 1000 and 5500 kilometers, bridging the gap between tactical SRBMs and strategic ICBMs. IRBMs can deliver nuclear or conventional warheads to targets across regional distances and are typically deployed on mobile launchers or in fixed silos. Examples include the US Pershing II and the Chinese DF-26, which have both been retired or are in active service. The IRBM category represents a significant strategic capability that fills the gap between theater and intercontinental strike ranges.}{catRM}

\dictentry{IRFNA}{Inhibited red fuming nitric acid, a storable oxidizer used in hypergolic propellant combinations with fuels like UDMH. IRFNA contains dissolved nitrogen tetroxide and a corrosion inhibitor that makes it compatible with standard metals at ambient temperature. It ignites spontaneously on contact with most organic fuels, providing reliable hypergolic ignition for storable liquid rocket engines. The inhibitor additive reduces the extremely corrosive nature of pure fuming nitric acid while preserving the oxidizer's hypergolic ignition characteristics and storage stability.}{catRM}

\dictentry{Iron Dome}{An Israeli short-range air defense system designed to intercept rockets, artillery, and mortar shells with ranges up to 70 kilometers. The system uses the Tamir interceptor guided by radar tracking and command control to destroy incoming projectiles with a proximity-fuzed warhead over populated areas. Iron Dome has achieved intercept rates exceeding 90 percent and has become the most combat-proven short-range air defense system in the world. The system's battle management computer selects which incoming threats require interception based on predicted impact points, conserving interceptors for threats that endanger protected areas.}{catRM}

\dictentry{Jet Tab}{A tab in the exhaust for steering that protrudes into the supersonic exhaust flow to create an asymmetric pressure distribution on one side of the nozzle. The resulting side force deflects the thrust vector without mechanical gimbaling of the entire engine. Jet tabs are simple and lightweight but cause some performance loss due to exhaust flow disturbance, and are used on some solid and liquid rocket motors for attitude control. The tab deflection angle and surface area are sized to provide the required control authority while minimizing the thrust loss penalty during operation.}{catRM}

\dictentry{Jet Vanes}{Vanes in the exhaust for steering that deflect part of the exhaust stream to produce a lateral force component for thrust vector control. Vanes are made of refractory materials such as graphite or tungsten to withstand the extreme temperatures of the exhaust flow. They are used on some tactical missiles and sounding rockets where simple and compact thrust vectoring is needed without the complexity of engine gimbaling. Jet vane erosion during operation progressively reduces control authority, which must be accounted for in the vehicle's control system design.}{catRM}

\dictentry{Kill Probability}{The probability of destroying the target, a statistical measure combining the probability of hit and the probability of kill given a hit. Kill probability depends on guidance accuracy, warhead lethality, target vulnerability, and the effectiveness of any countermeasures employed. Mission planners use kill probability to assess weapon effectiveness and determine the number of missiles required to achieve a desired confidence of target destruction. Kill probability analysis integrates engagement geometry, sensor performance, and weapon system reliability into a single metric for operational planning.}{catRM}

\dictentry{Kinetic Kill Vehicle}{A hit-to-kill interceptor that destroys the target by direct body-to-body impact at hypersonic closing speeds without using an explosive warhead. The KKV uses onboard sensors and divert thrusters to achieve the meter-level accuracy required for a direct collision in space. Systems such as the THAAD interceptor and the exoatmospheric kill vehicle on the GBI rely on kinetic kill technology for ballistic missile defense. The KKV must autonomously acquire, track, and maneuver to intercept the target within seconds of separating from its booster stage.}{catRM}

\dictentry{Laser Homing}{Homing using reflected laser energy from the target illuminated by the launch platform or a forward observer with a laser designator. The missile's seeker detects the reflected laser spot and steers to keep the target centered in its field of view. Laser-guided weapons offer high accuracy in clear weather but require line-of-sight illumination throughout the terminal phase, limiting use in poor visibility. The laser designator must maintain aim on the target until impact, which exposes the designator operator to risk and limits engagement geometry options.}{catRM}

\dictentry{Launch Azimuth}{The direction of launch measured from north, which determines the initial ground track and the resulting orbital plane for space launch vehicles. Launch azimuth is carefully chosen to achieve the desired orbital inclination while avoiding overflight of populated areas during the boost phase. For launches from Cape Canaveral, a typical azimuth of 28.5 degrees produces an orbit inclined at 28.5 degrees to the equator. Changes in launch azimuth directly affect the achievable orbital inclination and the range safety considerations during the ascent trajectory.}{catRM}

\dictentry{Launch Detection}{Detecting a missile launch by satellites or radar using infrared sensors on early-warning satellites or ground-based radar systems to identify the intense heat signature of a rocket motor exhaust plume. Space-based sensors such as the SBIRS constellation can detect launches within seconds, providing critical early warning to national command authorities. Detection data is immediately processed to classify the missile type, estimate its trajectory, and predict its impact point. The speed and accuracy of launch detection directly determines the available warning time for defensive response and civilian evacuation.}{catRM}

\dictentry{Launch Pad}{The facility from which a rocket launches, providing a stable platform, propellant loading infrastructure, electrical and data connections, and the flame trench and deflector for exhaust management. Launch pads include umbilical towers for pre-launch servicing, environmental control systems for payload conditioning, and lightning protection. Major launch complexes such as LC-39 at Kennedy Space Center and LC-1 at Baikonur Cosmodrome support operations for the largest launch vehicles. Pad turnaround time between launches is a key factor in launch campaign cadence and overall range productivity.}{catRM}

\dictentry{Launch Rail}{A rail guiding the rocket during initial ascent that provides directional constraint and stability until the rocket gains sufficient velocity for aerodynamic control. Launch rails are typically several meters long and aligned to the planned launch azimuth and elevation angle. They are used for tactical missiles, sounding rockets, and small launch vehicles where a fixed launch tower is impractical. The rail must be precisely aligned and rigidly supported to ensure the rocket exits at the correct angle without inducing excessive tip-off rates.}{catRM}

\dictentry{Launch Tower}{A structure supporting the rocket before launch that provides access platforms, propellant and electrical umbilicals, and environmental conditioning to the vehicle. The tower includes swing arms that retract at lift-off and may incorporate lightning protection masts and a crew access arm for crewed missions. Launch towers are typically destroyed or moved away from the vehicle on rails or pivots during the final seconds before liftoff. The tower design must accommodate the vehicle configuration, environmental conditions, and the servicing requirements of the specific launch campaign.}{catRM}

\dictentry{Launch Tube}{A tube for launching guided missiles that encloses the missile during the initial acceleration phase, protecting the launch platform and allowing launch from enclosed spaces such as submarine tubes or vehicle-mounted canisters. The launch tube provides structural support and guidance during the missile's initial motor ignition and exit. Tube-launched missiles are common in tactical applications where rapid deployment and platform protection are essential. The tube must withstand the thermal and pressure loads of missile motor ignition while providing a clear, unobstructed exit path for the missile.}{catRM}

\dictentry{Launch Vehicle}{A rocket used to launch payloads into space, consisting of one or more stages of propulsion, guidance, and structural systems designed to accelerate the payload to orbital velocity. Launch vehicles range from small expendable rockets for microsatellites to heavy-lift vehicles capable of delivering tens of thousands of kilograms to low Earth orbit. The selection of a launch vehicle depends on payload mass, target orbit, reliability requirements, and cost constraints. Launch vehicle development represents a multi-year, multi-billion-dollar investment that must balance performance, reliability, and commercial viability.}{catRM}

\dictentry{Lethality}{The damaging effect of a warhead on the target, determined by the warhead type, explosive yield, fragment characteristics, and the proximity and angle of detonation. Lethality is quantified by the probability of kill as a function of miss distance and is a key parameter in weapon effectiveness assessment. Warhead lethality can be enhanced through directed fragmentation, shaped charges, or proximity fuze optimization. Lethality modeling integrates weapon effects data with target vulnerability assessments to predict the expected damage for a given engagement scenario.}{catRM}

\dictentry{LH2}{Liquid hydrogen, a cryogenic rocket fuel stored at minus 253 degrees Celsius with the highest specific impulse of any practical chemical propellant when burned with liquid oxygen. LH2 has a very low density requiring large tank volumes, and its extreme cold demands insulated cryogenic storage and handling systems. Despite these challenges, LH2 is used on high-performance upper stages and main engines such as the SSME and RL10. The low density of LH2 means that tank volume is large relative to the propellant mass, requiring careful structural design to minimize the overall vehicle mass penalty.}{catRM}

\dictentry{Lift-Off}{The moment the rocket leaves the ground when engine thrust exceeds vehicle weight and the hold-down release mechanism allows the vehicle to begin ascent. Lift-off is the culmination of the countdown sequence and marks the transition from ground operations to flight. The first seconds after lift-off are the most critical as the vehicle accelerates slowly and is vulnerable to wind shear and asymmetric thrust conditions. Hold-down posts ensure that all engines reach full thrust and nominal performance before the vehicle is released to begin its ascent.}{catRM}

\dictentry{Lifting Trajectory}{A trajectory using aerodynamic lift to extend range or modify the flight path, typically employed by hypersonic glide vehicles and certain maneuvering reentry vehicles. Unlike a purely ballistic path, a lifting trajectory generates aerodynamic forces that allow cross-range and down-range maneuvering, making the vehicle's impact point unpredictable. This approach reduces peak deceleration and heating compared to a ballistic reentry but requires robust thermal protection and advanced guidance. The lift-to-drag ratio of the vehicle determines the achievable cross-range and the overall maneuvering capability during the glide phase.}{catRM}

\dictentry{Liquid Hydrogen}{A cryogenic fuel used in high-performance rockets, stored at minus 253 degrees Celsius as a transparent, low-density liquid. When burned with liquid oxygen, LH2 produces the highest specific impulse of any practical chemical propellant combination, making it the fuel of choice for upper-stage engines. Its low density requires large tank volumes, and its cryogenic nature demands vacuum-jacketed insulation and careful thermal management throughout ground and flight operations. LH2 boil-off management is a critical operational concern during ground handling and extended countdown sequences.}{catRM}

\dictentry{Liquid Injection Thrust Vectoring}{Injection of liquid into the nozzle for thrust vectoring, where a secondary fluid is injected at an angle into the divergent nozzle section to create a shock wave and asymmetric pressure distribution. The resulting lateral force deflects the thrust vector without mechanically moving the engine. LITV systems are simpler and lighter than mechanical gimbaling but produce some performance loss due to the injected fluid's momentum. The secondary injection fluid must be compatible with the exhaust environment and must not cause thermal damage to the nozzle wall at the injection point.}{catRM}

\dictentry{Liquid Oxygen}{A cryogenic oxidizer used with hydrogen and kerosene that is stored at minus 183 degrees Celsius as a pale blue liquid. LOX is the most common oxidizer in high-performance rocket engines due to its high density, excellent combustion performance, and relative non-toxicity. Handling LOX requires specialized cryogenic equipment and careful contamination control because liquid oxygen is strongly reactive with many organic materials. LOX conditioning and tanking operations require strict adherence to cleanliness procedures to prevent catastrophic reactions with hydrocarbon contamination.}{catRM}

\dictentry{Liquid Rocket Engine}{An engine using liquid fuel and oxidizer pumped to the combustion chamber where they are mixed and burned to produce high-velocity exhaust and thrust. Liquid engines offer throttleability, restart capability, and higher specific impulse compared to solid motors, but at the cost of greater system complexity. They are used on launch vehicle upper stages, main engines, and spacecraft propulsion where performance and controllability are paramount. The development of a liquid rocket engine typically requires years of design, testing, and qualification before flight certification is achieved.}{catRM}

\dictentry{LNG}{Liquefied natural gas as rocket fuel, primarily methane stored at minus 162 degrees Celsius, offering a middle ground between kerosene and liquid hydrogen in terms of density and specific impulse. LNG produces minimal coking and soot compared to kerosene, simplifying engine reuse, and is denser and easier to store than liquid hydrogen. Methane-fueled engines such as the SpaceX Raptor and Blue Origin BE-4 are being developed for next-generation reusable launch vehicles. LNG's favorable properties for reuse, including clean combustion and moderate cryogenic requirements, make it attractive for high-cadence operational launch systems.}{catRM}

\dictentry{Lofted Trajectory}{A higher than normal ballistic trajectory that trades range for reduced flight time and a steeper reentry angle. Lofted trajectories are used when the target is at shorter range than the missile's maximum capability or when a faster time of flight is desired to reduce the enemy's reaction time. The steeper reentry angle also reduces the effectiveness of some terminal-phase missile defense systems by limiting their engagement window. The tradeoff between flight time, reentry angle, and impact velocity must be optimized for each specific mission and target set.}{catRM}

\dictentry{LOX}{Liquid oxygen, a cryogenic oxidizer stored at minus 183 degrees Celsius that is the most widely used oxidizer in high-performance rocket engines. LOX is relatively inexpensive, non-toxic, and provides excellent combustion performance when paired with hydrogen, kerosene, or methane fuels. It must be handled with care because it strongly accelerates combustion and can cause spontaneous ignition of materials that would not normally burn in air. LOX system design includes strict material compatibility requirements and contamination control procedures to prevent dangerous reactions with organic materials.}{catRM}

\dictentry{Mach Diamond}{Diamond-shaped shock patterns in exhaust visible as a repeating series of luminous diamond shapes in the exhaust plume. Mach diamonds form when the exhaust exits the nozzle at a different pressure than the ambient atmosphere, creating oblique shock waves that reflect off the plume boundaries. The number and spacing of Mach diamonds provide visual information about the nozzle expansion ratio and the degree of overexpansion or underexpansion. These standing wave patterns are a direct visual indicator of the pressure mismatch between the exhaust jet and the surrounding atmosphere.}{catRM}

\dictentry{Main Combustion Chamber}{The primary chamber where full combustion occurs at the designed chamber pressure and temperature, producing the majority of the engine's thrust. In staged combustion engines, the main chamber receives fully mixed and preheated propellants from the preburners through the injector. The main chamber wall is typically regeneratively cooled by circulating one of the propellants before it enters the injector. The chamber design must withstand extreme thermal and pressure loads while maintaining structural integrity throughout the engine burn duration.}{catRM}

\dictentry{Main Propellant Valve}{The valve controlling propellant flow to the engine that governs engine start, shutdown, and throttling by regulating the flow of fuel and oxidizer to the injector. Main propellant valves must operate reliably at cryogenic or elevated temperatures and respond quickly to start and stop commands. They are critical for safe engine operation and are designed with redundancy and leak-tight sealing to prevent uncommanded propellant flow. The valve opening and closing profiles are carefully controlled to manage engine start and shutdown transients and prevent dangerous combustion instabilities.}{catRM}

\dictentry{Mass Flow Rate}{The mass of propellant flowing per unit time through the engine, typically expressed in kilograms per second. Mass flow rate determines the thrust level for a given exhaust velocity and directly affects chamber pressure, combustion temperature, and nozzle performance. Accurate mass flow rate measurement and control are essential for engine health monitoring and performance prediction during flight. The ratio of fuel to oxidizer mass flow rates, known as the mixture ratio, is a critical parameter that must be precisely controlled for optimal combustion efficiency.}{catRM}

\dictentry{Max Q}{The point of maximum dynamic pressure during ascent, the moment when the combination of increasing vehicle velocity and decreasing atmospheric density produces the highest aerodynamic pressure on the vehicle structure. Max Q typically occurs between 60 and 90 seconds after launch at an altitude of 10 to 15 kilometers. Vehicle structures are designed to withstand the loads at max Q, and some engines reduce thrust temporarily to stay within structural limits. The dynamic pressure at max Q determines the maximum aerodynamic bending moments and axial loads experienced by the vehicle structure during ascent.}{catRM}

\dictentry{Maximum Range Trajectory}{The trajectory achieving the greatest range for a given missile and propellant load, balancing the tradeoff between flight time, drag losses, and gravitational losses. Maximum range is achieved by optimizing the launch angle, pitch program, and throttle profile to maximize downrange distance. For ballistic missiles, the maximum range trajectory is typically a 45-degree launch angle in a vacuum, but atmospheric effects and staging modify this optimum. Range optimization must account for the Earth's rotation, atmospheric density profile, and the vehicle's thrust-to-weight characteristics throughout the burn.}{catRM}

\dictentry{MECO}{Main engine cutoff, the moment when the primary propulsion system shuts down after completing its planned burn. MECO is typically triggered by achieving a target velocity, propellant depletion, or a commanded shutdown based on trajectory requirements. The precise timing of MECO directly affects the achieved orbit or trajectory, and post-MECO events such as stage separation and second-stage ignition follow in a tightly sequenced timeline. MECO sensor redundancy ensures that the shutdown command is executed reliably even if primary sensors provide conflicting data.}{catRM}

\dictentry{Metal Nozzle}{A nozzle made of high-temperature metal such as niobium alloy, molybdenum, or rhenium that can withstand combustion temperatures through regenerative cooling or radiative cooling. Metal nozzles are reusable and offer good resistance to thermal shock, making them suitable for restartable upper-stage engines. Refractory metal nozzles are often used in pressure-fed spacecraft thrusters where simplicity and reliability are more important than maximum performance. The choice of refractory metal depends on the operating temperature, burn duration, and the required number of thermal cycles over the engine's operational life.}{catRM}

\dictentry{Meteorological Rocket}{A rocket for atmospheric measurements that carries scientific instruments to measure temperature, pressure, wind speed, and composition of the upper atmosphere. Met rockets are typically small solid-propellant sounding rockets launched to altitudes of 50 to 100 kilometers, providing data that cannot be obtained by balloons or satellites. The data collected is used for weather forecasting, atmospheric research, and validating climate models. Meteorological rockets deploy instrumented payloads on parachutes during descent to obtain vertical profiles of atmospheric parameters throughout the mesosphere and lower thermosphere.}{catRM}

\dictentry{Midcourse Guidance}{The intermediate phase of missile guidance between boost and terminal phases, where the missile corrects its trajectory using inertial updates, GPS fixes, or external command data. Midcourse guidance compensates for launch errors and steering inaccuracies accumulated during boost and aligns the missile for an optimal terminal approach. For ballistic missile defense, midcourse guidance is critical for positioning the interceptor at the predicted intercept point before the terminal seeker activates. The midcourse phase represents the longest duration of most missile flights and offers the greatest opportunity for trajectory correction.}{catRM}

\dictentry{Minimum Energy Trajectory}{The most efficient trajectory for a given range that minimizes the total propellant required by optimizing the launch angle and flight path. The minimum energy trajectory balances the tradeoff between gravitational losses from flying too steep and drag losses from flying too shallow. For intercontinental range, this corresponds to a launch angle near 45 degrees and a trajectory that spends most of its time above the sensible atmosphere. Deviations from the minimum energy trajectory are intentional when other mission objectives such as reduced flight time or defense penetration are prioritized.}{catRM}

\dictentry{MIRV}{Multiple independently targetable reentry vehicles, a technology allowing a single ballistic missile to carry several warheads each aimed at a different target. After boost and midcourse, the post-boost vehicle dispenses individual reentry vehicles on separate trajectories to their respective aim points. MIRV dramatically increases the destructive potential of each missile and complicates missile defense by requiring interceptors to engage multiple targets from a single launch. The post-boost vehicle must precisely position and release each reentry vehicle with accurate attitude and velocity conditions for independent targeting.}{catRM}

\dictentry{Missile}{A guided weapon system consisting of a propulsion section, guidance and navigation system, warhead or payload, and control surfaces that steer it toward a predetermined or detected target. Missiles can be launched from ground, sea, air, or subsurface platforms and range from short-range tactical weapons to intercontinental strategic systems. The term distinguishes missiles from unguided rockets by their ability to autonomously or semi-autonomously correct their trajectory during flight. Modern missiles integrate multiple guidance modes, advanced warheads, and sophisticated counter-countermeasure capabilities to defeat increasingly capable defensive systems.}{catRM}

\dictentry{MMH}{Monomethylhydrazine, a hypergolic fuel with the chemical formula CH3NHNH2 commonly used in bipropellant rocket engines with nitrogen tetroxide as the oxidizer. MMH offers good performance, moderate freezing point, and long-term storability, making it the preferred fuel for spacecraft thrusters and orbital maneuvering engines. It is toxic and must be handled with extreme care using specialized personal protective equipment and closed-loop propellant transfer systems. MMH's favorable properties for long-duration space missions include excellent thermal stability and compatibility with common spacecraft materials over multi-year storage periods.}{catRM}

\dictentry{Mobile Launcher}{A transportable launch platform that can be moved by road, rail, or air to pre-surveyed launch positions, providing operational flexibility and survivability against preemptive strike. Mobile launchers for ICBMs such as the Russian Topol-M and Chinese DF-41 can disperse across a wide area and launch from unprepared sites. For tactical missiles, mobile launchers enable rapid repositioning and shoot-and-scoot tactics that complicate enemy targeting. The launcher must provide a stable firing platform while remaining sufficiently compact and lightweight for rapid transport over road and rail networks.}{catRM}

\dictentry{Monopropellant}{A single propellant decomposed by a catalyst to produce hot gas for thrust, most commonly hydrazine passed over an iridium or ruthenium catalyst bed. Monopropellant thrusters are the simplest form of liquid propulsion, requiring no oxidizer, ignition system, or mixing hardware. They are widely used for spacecraft attitude control, station-keeping, and orbital maneuvering where moderate thrust and high reliability are more important than maximum performance. The catalyst bed life and degradation rate are critical factors in determining the operational lifetime of monopropellant thruster systems.}{catRM}

\dictentry{Motor Case}{The structural casing of a solid rocket motor that contains the propellant grain and withstands the high internal pressures generated during combustion. Motor cases are manufactured from maraging steel, titanium alloys, or filament-wound composite materials such as carbon fiber epoxy to achieve the best strength-to-weight ratio. The case design must account for thermal loads, handling loads, and the pressure loads from both propellant weight and combustion pressure. Composite motor cases offer significant weight savings over metallic designs but require careful fiber placement and cure process control to achieve consistent structural performance.}{catRM}

\dictentry{Multi-Stage Rocket}{A rocket with multiple stages discarded during ascent, where each stage contains its own engines and propellant and is jettisoned after exhausting its fuel to reduce the dead weight carried by subsequent stages. Staging dramatically improves payload capability by allowing each stage to operate at a favorable mass ratio. The Saturn V used three stages and the Falcon 9 uses two stages, with each separation event carefully sequenced to minimize performance loss. The optimal number of stages, propellant loading, and separation altitudes are determined through trajectory optimization analysis for each specific mission.}{catRM}

\dictentry{Neutral Burning}{Burning with constant surface area that produces a steady thrust level throughout the burn. Neutral burning grains are typically cylindrical with perforations sized so that the increasing surface of the burning perforation offsets the decreasing outer surface. This configuration is desirable for applications requiring constant thrust, such as main propulsion motors and sustainers. Achieving truly neutral burning requires precise control of the grain geometry and perforation dimensions during the casting and curing process.}{catRM}

\dictentry{Nitrogen Pressurization}{Pressurization using nitrogen gas to maintain tank pressure and force propellant to the engines, typically used on smaller and less performance-critical systems where the simplicity outweighs the weight penalty compared to helium. Nitrogen is stored in high-pressure vessels and regulated to the required tank pressure before introduction above the propellant. This approach is common on tactical missile systems and ground-launched rockets where operational simplicity is valued. Nitrogen pressurization eliminates the need for complex helium pressurization hardware while providing adequate tank pressurization for short-burn tactical applications.}{catRM}

\dictentry{Nitrogen Tetroxide}{A storable oxidizer for hypergolic propellants with the chemical formula N2O4 that exists as a reddish-brown liquid at room temperature. N2O4 dissociates into NO2 at elevated temperatures and is hypergolic with most organic fuels including UDMH, MMH, and hydrazine. Its storability at ambient temperature and reliable spontaneous ignition make it the standard oxidizer for storable bipropellant spacecraft and missile propulsion systems. N2O4 vapor is toxic and corrosive, requiring specialized handling equipment, protective gear, and containment systems during ground operations.}{catRM}

\dictentry{Nozzle}{A duct accelerating exhaust gases to produce thrust by converting the thermal energy of high-pressure, high-temperature combustion products into directed kinetic energy. The nozzle converges to a sonic throat and then diverges to accelerate the flow to supersonic velocities, with the expansion ratio determining the final exhaust velocity. Nozzle design is a critical element of rocket engine performance, affecting thrust, efficiency, and operational envelope. The nozzle contour, throat area, and exit diameter must be optimized together to achieve the desired performance across the intended operating altitude range.}{catRM}

\dictentry{Nozzle Closure}{A seal protecting the nozzle before ignition that prevents moisture, debris, and foreign objects from entering the combustion chamber and nozzle during storage, transport, and handling. Nozzle closures are typically thin metal or composite diaphragms that are ruptured by chamber pressure during motor ignition. They must withstand environmental loads and maintain a hermetic seal throughout the motor's service life. The closure burst pressure must be carefully calibrated to ensure reliable rupture at motor ignition without premature failure during handling or thermal cycling.}{catRM}

\dictentry{Nozzle Contour}{The shape of a rocket nozzle wall that determines the expansion process and ultimately the thrust produced by the engine. Optimal contours such as the Rao minimum-length bell nozzle maximize thrust for a given length by rapidly expanding the flow near the throat and gradually turning it back to axial direction. The contour shape directly affects performance, flow separation characteristics, and the thermal loads on the nozzle wall. Computational fluid dynamics optimization is used to refine the nozzle contour for specific mission requirements and operating conditions.}{catRM}

\dictentry{Nozzle Extension}{An extendable nozzle section for altitude compensation that increases the expansion ratio after the vehicle leaves the dense atmosphere. Deployable nozzle extensions are stowed during ground-level operation to prevent flow separation and then extended in space to improve vacuum performance. This approach is used on upper-stage engines like the RL10 and Vinci to optimize sea-level and vacuum operation within a single engine design. The deployment mechanism must reliably extend and lock the nozzle extension in the zero-gravity vacuum environment of space.}{catRM}

\dictentry{Nozzle Throat Insert}{A high-temperature material at the nozzle throat that resists the extreme thermal and erosive environment where gas velocity reaches sonic conditions and heat flux is highest. Throat inserts are made from graphite, carbon-carbon composite, tungsten, or rhenium to minimize erosion and maintain a constant throat area throughout the burn. Throat erosion directly reduces chamber pressure and degrades engine performance, making insert material selection critical for long-duration burns. The insert must be securely bonded to the surrounding nozzle structure to prevent gas leakage and maintain the designed throat geometry.}{catRM}

\dictentry{Nuclear Warhead}{A nuclear explosive warhead that derives its destructive energy from nuclear fission, fusion, or a combination of both, producing blast, thermal radiation, and ionizing radiation effects far exceeding conventional explosives. Nuclear warheads are carried by ICBMs, SLBMs, and some IRBMs as the primary strategic deterrent force of nuclear-armed nations. Modern thermonuclear warheads use a fission primary to ignite a fusion secondary, achieving yields measured in hundreds of kilotons. Warhead reliability and safety are maintained through rigorous quality assurance, surveillance programs, and certified maintenance procedures throughout the weapon's operational life.}{catRM}

\dictentry{Optimal Expansion}{Nozzle expansion matching ambient pressure, where the exhaust exits the nozzle at the same pressure as the surrounding atmosphere to maximize thrust efficiency. At optimal expansion, there are no pressure losses at the nozzle exit plane, and all available energy is converted to directed kinetic energy. Since ambient pressure changes with altitude, a fixed-geometry nozzle can only be optimally expanded at one altitude, which is why altitude-compensating designs are attractive. Achieving optimal expansion across the full flight envelope remains one of the key challenges in nozzle design for single-stage-to-orbit and multi-mission vehicles.}{catRM}

\dictentry{Over-the-Horizon Radar}{Radar seeing beyond the horizon by reflecting high-frequency radio waves off the ionosphere to detect targets at ranges of hundreds to thousands of kilometers. OTH radars can detect incoming ballistic missiles, aircraft, and ships well beyond the line-of-sight limitation of conventional microwave radar. They provide early warning and cueing for missile defense systems, though with lower accuracy than direct-line-of-sight radars. OTH radar performance depends on ionospheric conditions, which vary with solar activity, time of day, and season, requiring adaptive frequency management.}{catRM}

\dictentry{Overexpansion}{Nozzle exit pressure below ambient, occurring when a nozzle designed for high-altitude operation operates at lower altitude where atmospheric pressure is higher than the exhaust pressure. Overexpansion reduces thrust efficiency and can cause flow separation from the nozzle wall, producing lateral loads and potential structural damage. Nozzle design must account for overexpansion conditions during the initial phase of sea-level launch. Controlled overexpansion is tolerable for short durations, but sustained operation in severely overexpanded conditions can cause destructive side loads and nozzle failure.}{catRM}

\dictentry{Oxidizer}{The substance supplying oxygen for combustion in a rocket, which reacts with the fuel to produce the high-temperature, high-pressure gases that generate thrust. Common oxidizers include liquid oxygen, nitrogen tetroxide, and hydrogen peroxide, each offering different tradeoffs in performance, storability, and toxicity. The choice of oxidizer is one of the most fundamental decisions in rocket propulsion design, affecting engine cycle, materials, and ground handling. Oxidizer selection influences every aspect of the propulsion system from injector material compatibility to tank design and ground support equipment requirements.}{catRM}

\dictentry{Passive Radar Homing}{Homing using emissions from the target where the missile's seeker detects and tracks radio frequency energy radiated by the target's own radar, communications, or electronic systems. The missile requires no illumination from the launch platform and gives the target no warning of the incoming threat until impact. Passive radar homing is particularly effective against naval vessels and airborne early warning aircraft that emit powerful radar signals. The seeker must be able to discriminate the target's emissions from background clutter and intentional jamming signals in a complex electromagnetic environment.}{catRM}

\dictentry{Patriot}{A US surface-to-air missile system providing medium-range air defense against aircraft, cruise missiles, and tactical ballistic missiles using the PAC-3 interceptor with hit-to-kill technology. The Patriot system consists of a phased-array radar, engagement control station, and launchers that can be positioned dispersed from the radar. Patriot has been extensively combat-tested in Gulf War, Iraq, and Ukraine operations, evolving through multiple upgrades to counter increasingly sophisticated threats. The PAC-3 variant represents a fundamental shift to kinetic kill technology, providing improved effectiveness against ballistic missile threats compared to the older blast-fragmentation warhead interceptors.}{catRM}

\dictentry{Payload Fairing}{The nose cone protecting the payload from aerodynamic pressure, heating, and acoustic loads during atmospheric ascent. The fairing is typically made of composite materials such as honeycomb sandwich panels with acoustic blanket lining and is jettisoned once the vehicle is above the sensible atmosphere. Fairing design must balance thermal protection, structural integrity, acoustic attenuation, and minimal weight to maximize payload delivery capability. Fairing acoustics are a significant design driver because the enclosed environment amplifies engine and aerodynamic noise that can damage sensitive payload components.}{catRM}

\dictentry{Payload Mass Fraction}{The ratio of payload mass to total mass at launch, representing the fraction of the vehicle's initial mass that is useful payload delivered to orbit. Payload mass fraction is the ultimate measure of launch vehicle efficiency and is typically only 2 to 4 percent for expendable vehicles. Improving payload mass fraction through structural weight reduction, better propellants, and optimized trajectories is a primary goal of launch vehicle design. Every incremental improvement in payload mass fraction translates directly into either increased revenue from additional payload capacity or reduced launch cost per kilogram delivered to orbit.}{catRM}

\dictentry{Penetration Aid}{A device helping a missile penetrate defenses by creating false targets, radar clutter, or chaff clouds that confuse and saturate the defensive system's sensors and interceptors. Penetration aids include decoys, jammers, chaff, and maneuvering reentry vehicles that complicate the discrimination problem for missile defense. They are essential for ensuring that at least some warheads reach their targets through an adversary's defensive shield. The selection and combination of penetration aids is tailored to the specific defensive system being countered and the threat environment expected during the engagement.}{catRM}

\dictentry{Phased Array Radar}{A radar with electronic beam steering that uses an array of small antenna elements whose individual phase is controlled to steer the radar beam without physically moving the antenna. Phased array radars can switch beam direction in microseconds, allowing rapid tracking of multiple targets simultaneously while performing search functions. Aegis SPY radars, Patriot radars, and THAAD radars all use phased array technology for multi-target engagement capability. The electronic beam steering enables near-simultaneous search, track, and missile guidance functions that would require multiple mechanically scanned radars.}{catRM}

\dictentry{Pintle Injector}{An injector with a central pintle for variable flow that produces a hollow-cone spray pattern by injecting one propellant axially through a central post and the other radially through an annular gap. The pintle design allows deep throttle capability because the injection velocity and mixing characteristics remain favorable over a wide flow range. The SpaceX Merlin engine uses a pintle injector, enabling the throttling required for propulsive landing. The pintle geometry provides inherent combustion stability across a wide operating range, which is particularly valuable for deeply throttled engine operations.}{catRM}

\dictentry{Plug Nozzle}{An altitude-compensating nozzle design featuring a central plug around which the exhaust flows, with ambient pressure acting on the external plume boundary to automatically adjust the effective expansion ratio. At sea level the atmospheric pressure confines the plume against the plug, while at altitude the plume expands further, adapting to ambient conditions without mechanical changes. Plug nozzles offer theoretical performance benefits across the flight envelope but present manufacturing and cooling challenges. The plug surface must withstand extreme heat flux while maintaining the contour accuracy needed for efficient gas expansion.}{catRM}

\dictentry{Pneumatic Valve}{A valve actuated by compressed gas, typically helium or nitrogen, that uses a gas piston or diaphragm to open or close the valve mechanism. Pneumatic actuators offer fast response times and reliable operation in the extreme thermal and vibration environment of rocket engines. They are commonly used for propellant isolation, sequencing, and emergency shutdown valves in both liquid and solid propulsion systems. Pneumatic valve design must ensure positive seating and leak-tight closure under the full range of operating pressures and temperatures encountered during engine operation.}{catRM}

\dictentry{Pogo Oscillation}{Longitudinal oscillation from engine-vehicle interaction where propellant flow pulsations couple with the structural vibration modes of the rocket, creating a self-excited oscillation along the vehicle's axis. Pogo can cause severe structural loads, reduced payload life, and even vehicle breakup if not suppressed. Pogo suppressors, which are gas-filled accumulators in the propellant lines, are installed to decouple the engine from the vehicle structure and damp the oscillation. The pogo phenomenon was first identified on the Saturn V program and remains a critical design consideration for all large multi-stage launch vehicles.}{catRM}

\dictentry{Poppet Valve}{A valve with a disc and seat where a disc-shaped poppet moves perpendicular to the seat to open or close the flow path. Poppet valves provide positive shutoff with low leakage and are commonly used as check valves, relief valves, and isolation valves in propellant feed systems. Their simple design and reliable sealing make them one of the most widely used valve types in rocket propulsion. Poppet valve seating force and response time are critical parameters that must be optimized for the specific flow and pressure conditions of each application.}{catRM}

\dictentry{Preburner}{A small chamber generating gas to drive the turbine by partially burning propellants in either a fuel-rich or oxidizer-rich mode to produce warm, high-pressure gas. The preburner operates at very high pressure and temperature, and its output drives the turbopump before the gas enters the main combustion chamber. Preburners are central to staged combustion cycle engines, enabling very high chamber pressures and engine performance. The preburner must produce stable, homogeneous combustion to avoid turbine blade damage from temperature non-uniformities or combustion-driven oscillations.}{catRM}

\dictentry{Pressure Fed Cycle}{An engine cycle using tank pressure for propellant feed where pressurized gas pushes the propellants from their tanks to the combustion chamber without turbopumps. This cycle offers maximum simplicity and reliability at the cost of heavier tanks that must withstand the pressurization loads. Pressure-fed engines are used on spacecraft thrusters and some upper stages where high reliability and simplicity outweigh the performance penalty of heavier tanks. The maximum achievable chamber pressure is limited by the tank pressure capability, which restricts pressure-fed engines to relatively low thrust and performance levels.}{catRM}

\dictentry{Pressure Regulator}{A device maintaining constant pressure by reducing high-pressure supply gas to a stable, lower output pressure for tank pressurization or pneumatic actuation. Pressure regulators must respond quickly to changing flow demands while maintaining precise output pressure regulation across a wide range of inlet conditions. They are critical components in propellant feed systems, controlling the pressure that drives propellant flow to the engine. Regulator stability and response characteristics directly affect tank pressurization uniformity and engine start and shutdown transient behavior.}{catRM}

\dictentry{Progressive Burning}{Burning with increasing surface area that produces rising thrust over the course of the burn. Progressive grain geometries such as star or finocyl configurations with appropriate dimensions increase the burning surface as the combustion front progresses. This configuration is useful for applications requiring increasing thrust, such as launch vehicle boosters that need maximum thrust as the vehicle becomes lighter from propellant consumption. The rate of thrust increase is controlled by the grain geometry and must be matched to the vehicle's structural limits and trajectory requirements.}{catRM}

\dictentry{Propellant}{The fuel and oxidizer mixture burned in a rocket engine, selected for its energy density, storability, handling characteristics, and compatibility with engine materials. Propellant choices range from solid composite charges to cryogenic liquid pairs like liquid oxygen and liquid hydrogen, each with distinct tradeoffs in performance, safety, and operational complexity. The propellant selection is one of the most fundamental decisions in rocket design, influencing virtually every other subsystem. Propellant properties determine the engine cycle, tank materials, ground handling procedures, and the overall logistics footprint of the launch system.}{catRM}

\dictentry{Propellant Line}{A line carrying propellant from the tanks to the engine, designed to withstand internal pressure, thermal expansion, vibration, and external loads while maintaining leak-tight integrity. Propellant lines are typically made from stainless steel, Inconel, or titanium alloys and are routed with flex joints and supports to accommodate vehicle structural deflection. Line sizing balances flow velocity against pressure drop and weight to ensure adequate propellant delivery under all operating conditions. Propellant line routing must avoid sharp bends and dead volumes that could trap gas or cause flow instabilities during engine operation.}{catRM}

\dictentry{Propellant Liner}{A liner bonding the grain to the case that provides an adhesive and insulating interface between the solid propellant and the motor case wall. The liner must withstand the thermal and pressure loads of motor operation while maintaining its bond to both the case and the grain. Liner materials are typically elastomeric compounds that are applied to the case interior before propellant casting and cure. The liner thickness and application quality directly affect motor safety, as liner failure can expose the case wall to direct combustion gases and cause catastrophic burnthrough.}{catRM}

\dictentry{Propellant Management}{The system managing propellant in tanks, including baffles, pickups, surface tension screens, and sumps that ensure gas-free propellant delivery to the engine regardless of vehicle orientation or acceleration state. In microgravity environments, propellant management devices are essential for directing liquid to the tank outlet since buoyancy does not separate gas from liquid. Proper PMD design prevents gas ingestion that could cause engine cavitation, instability, or shutdown. Propellant management devices must function reliably across the full range of acceleration environments from zero-g coast to high-g boost phases.}{catRM}

\dictentry{Propellant Mass Fraction}{The ratio of propellant mass to total mass, representing the fraction of the vehicle's launch weight that is propellant available for generating thrust. Higher propellant mass fractions mean more delta-v for a given vehicle size and are achieved through lightweight structures and efficient tank design. Advanced composite cases and optimized grain designs push propellant mass fractions above 90 percent in modern solid rocket motors. The propellant mass fraction is a key figure of merit that directly determines the vehicle's velocity capability and payload delivery performance.}{catRM}

\dictentry{Propellant Settling}{Using small thrusters to settle propellant by creating a slight acceleration that pushes liquid propellant to the tank outlet, ensuring gas-free priming of the main engine before ignition. In microgravity or during coast phases, propellant floats randomly in the tank and may not cover the outlet. Settling burns are typically performed seconds before main engine start and use dedicated small thrusters or reaction control jets. The settling acceleration must be sufficient to overcome surface tension effects and ensure complete liquid coverage of the tank outlet.}{catRM}

\dictentry{Propellant Tank}{A tank storing rocket propellant that must withstand pressurization loads, thermal environments, and dynamic loads while minimizing structural weight. Tanks are typically cylindrical with hemispherical domes for optimal pressure vessel efficiency and are made from aluminum, stainless steel, or composite materials. Internal features such as anti-slosh baffles, propellant management devices, and pressurization interfaces are integrated into the tank design. Tank design must account for thermal contraction when cryogenic propellants are loaded and for structural loads during all phases of flight and ground handling.}{catRM}

\dictentry{Proximity Fuse}{A fuse detonating near the target using radar, laser, or infrared sensors to measure the distance to the target and trigger the warhead at the optimal standoff for maximum lethality. Proximity fuzing greatly increases the probability of kill compared to contact-only fuzing because it does not require a direct hit. Modern multi-mode fuzes can select between proximity, contact, and delay modes depending on the target type and engagement geometry. The fuse must discriminate between the target and background clutter, chaff, and countermeasures while operating reliably at extreme closing velocities.}{catRM}

\dictentry{Pyro Valve}{A valve actuated by pyrotechnic charge that uses a small explosive or thermal cartridge to drive a piston or rupture a diaphragm, providing a one-shot but extremely reliable actuation. Pyro valves are used for propellant isolation, emergency venting, and flight termination applications where absolute reliability and fast response outweigh the inability to reopen. They are common on solid rocket motors and spacecraft propulsion systems for safety-critical functions. Pyro valve actuation time is typically measured in milliseconds, providing nearly instantaneous response for time-critical safety and sequencing functions.}{catRM}

\dictentry{Pyrogen Igniter}{A pyrotechnic igniter generating hot gas that is used to initiate combustion in the main chamber of a liquid rocket engine. The pyrogen produces a stream of hot, particle-laden gas that impinges on the propellant in the chamber, raising it to ignition temperature. Pyrogen igniters are reliable and capable of providing the large energy input needed to start large liquid engines such as the F-1 and RD-180. The pyrogen itself requires a separate ignition source, typically an electrical squib, which initiates the pyrogen charge to produce the main ignition energy pulse.}{catRM}

\dictentry{Pyrotechnic Igniter}{An igniter using pyrotechnic compounds such as boron-potassium nitrate or zirconium-nickel to generate intense heat and flame for initiating combustion. Pyrotechnic igniters are initiated by an electrical squib and provide a short-duration, high-intensity energy pulse. They are widely used for solid rocket motor ignition, cartridge-start gas generators, and emergency pyrotechnic devices throughout missile and launch vehicle systems. The pyrotechnic composition must be carefully formulated to deliver the required energy output while maintaining shelf stability and safe handling characteristics throughout the weapon system's service life.}{catRM}

\dictentry{Radiative Cooling}{Cooling by thermal radiation from the nozzle, where the nozzle wall emits infrared radiation to the environment to reject heat absorbed from combustion gases. Radiative cooling is effective for nozzle extensions and lightweight engines operating at moderate heat flux, and requires no coolant plumbing or pump power. The nozzle material must withstand the elevated wall temperatures, typically using refractory metals or carbon composites with high emissivity coatings. Radiative cooling efficiency depends on the wall temperature, surface emissivity, and the view factor between the nozzle wall and the surrounding environment.}{catRM}

\dictentry{Ramjet Missile}{A missile using a ramjet engine that compresses incoming air by the vehicle's forward speed through an inlet diffuser, mixes it with fuel, and burns it in a continuous combustion process to produce thrust. Ramjets cannot start from rest and require a rocket booster or drop launch to achieve sufficient speed for air intake compression. They offer higher sustained speeds and longer range than rocket-powered missiles of similar size for atmospheric flight profiles. Ramjet efficiency increases with speed, making them most effective at supersonic velocities where the ram compression ratio is highest.}{catRM}

\dictentry{Range Safety}{Protecting the public from launch hazards through flight termination systems, restricted launch corridors, and real-time tracking that ensures a deviating vehicle is destroyed before reaching populated areas. Range safety officers monitor every launch from designated tracking stations and have the authority and capability to command vehicle destruction at any point during flight. Safety protocols include exclusion zones on the ground and in airspace that are cleared before every launch event. Modern range safety systems increasingly use autonomous flight termination capability to reduce dependence on ground-based tracking and command links.}{catRM}

\dictentry{Rao Nozzle}{A nozzle contoured for maximum thrust using the method of characteristics to design a minimum-length bell shape that achieves near-optimal expansion. The Rao contour turns the exhaust flow from the throat to axial direction in the shortest possible distance, minimizing nozzle weight while maximizing performance. This nozzle design is the standard for modern rocket engines and is derived from gas dynamics optimization theory. The Rao bell contour typically achieves 98 to 99 percent of the theoretical maximum thrust coefficient for a given expansion ratio and length.}{catRM}

\dictentry{Reaction Control System}{Small thrusters for attitude control in space that provide torque for pitch, yaw, and roll corrections as well as small translation maneuvers. RCS thrusters typically use monopropellant hydrazine or bipropellant combinations and are positioned around the vehicle to produce the required control moments. They are essential for maintaining spacecraft orientation, performing orbital corrections, and executing proximity operations such as docking. RCS propellant consumption must be carefully budgeted over the mission lifetime, as the propellant supply is finite and cannot be replenished in orbit.}{catRM}

\dictentry{Reentry Phase}{The final phase returning through the atmosphere when the reentry vehicle encounters increasing aerodynamic heating and deceleration as it descends from space to the target. During reentry, the vehicle experiences peak heating rates, maximum deceleration, and communication blackout caused by the ionized plasma sheath. Thermal protection systems such as ablative heat shields or ceramic tiles protect the vehicle structure from temperatures exceeding 10000 Kelvin at the stagnation point. The reentry trajectory must be carefully controlled to manage heating rates, deceleration loads, and landing accuracy within the design limits of the thermal protection system.}{catRM}

\dictentry{Reentry Vehicle}{The part of a missile carrying the warhead through reentry, designed with a heat shield and aerodynamic shape to survive the extreme thermal and mechanical loads of atmospheric reentry. RVs are typically small, dense, and spin-stabilized to maintain a predictable trajectory from the upper atmosphere to the target. MIRV systems deploy multiple independent reentry vehicles from a single post-boost vehicle, each aimed at a different target. The RV shape, mass distribution, and center of gravity location are carefully controlled to ensure predictable aerodynamic behavior and accurate impact point performance.}{catRM}

\dictentry{Regenerative Cooling}{Cooling a nozzle by circulating propellant through channels in the nozzle or chamber wall before the propellant enters the injector for combustion. The propellant absorbs heat from the wall and is preheated in the process, recovering energy that would otherwise be lost. Regenerative cooling is the primary thermal management method for high-performance liquid rocket engines and enables sustained operation at very high chamber pressures and temperatures. The cooling channel geometry, wall thickness, and propellant flow distribution must be optimized to prevent local hot spots that could cause wall burnthrough.}{catRM}

\dictentry{Regressive Burning}{Burning with decreasing surface area that produces declining thrust over the course of the burn. Regressive burning is achieved with grain geometries such as end-burning cylinders or finocyl shapes where the combustion front moves inward, reducing the exposed surface. This configuration is useful for upper-stage motors or sustainers where decreasing thrust as the vehicle becomes lighter maintains a favorable thrust-to-weight ratio. The rate of thrust regression is controlled by the grain geometry and must be precisely predictable for accurate trajectory performance.}{catRM}

\dictentry{Relief Valve}{A valve preventing overpressure by automatically opening and venting gas or propellant when the system pressure exceeds a predetermined safe limit. Relief valves protect tanks, lines, and components from catastrophic failure due to overpressurization caused by thermal expansion, regulator malfunction, or pyrotechnic events. They are calibrated to open at a specific pressure and reclose when pressure returns to a safe level. Relief valve sizing must account for the maximum expected relief flow rate to prevent dangerous pressure buildup during overpressure events.}{catRM}

\dictentry{Resonator}{A device damping specific frequencies of pressure oscillation in the combustion chamber by providing a tuned cavity or volume that absorbs acoustic energy at the problem frequencies. Resonators work on the Helmholtz principle, where the cavity volume and neck geometry determine the resonant frequency of energy absorption. They are used as supplements to baffles and injector design changes to address combustion instability during engine development. Resonator tuning requires detailed acoustic analysis of the chamber modes and iterative hot-fire testing to verify effective suppression of the target instability frequencies.}{catRM}

\dictentry{Retro Rocket}{A rocket firing opposite to the direction of travel to decelerate the vehicle for landing, orbit reduction, or reentry initiation. Retro rockets are used for propulsive landing of reusable boosters, deorbiting spacecraft, and slowing landers on planetary surfaces. The thrust must be precisely timed and controlled to achieve the desired velocity change without tumbling or uncontrolled rotation. Retro rocket ignition timing and thrust level must be precisely calculated to achieve the target velocity decrement while accounting for the vehicle's changing mass and orientation during the deceleration burn.}{catRM}

\dictentry{Reusable Launch Vehicle}{A launch vehicle with recoverable stages designed to be refurbished and reflown, dramatically reducing the per-launch cost compared to expendable vehicles. Reusability requires robust structures, restartable engines, and thermal protection systems capable of withstanding multiple entry and landing cycles. The SpaceX Falcon 9 first stage and the developing Starship system represent the leading examples of operational reusable launch technology. Achieving economically viable reusability requires minimizing refurbishment time and cost while maintaining the structural and system reliability needed for safe repeated flight operations.}{catRM}

\dictentry{Rocket Engine}{A thrust chamber where propellant is burned or heated and expelled at high velocity to produce thrust through Newton's third law. Rocket engines carry both fuel and oxidizer, enabling them to operate in the vacuum of space as well as within the atmosphere. They range from small attitude control thrusters producing millinewtons of force to massive main engines generating millions of newtons for launch vehicle boost. Rocket engine development requires extensive testing and qualification to demonstrate the reliability and performance needed for flight certification.}{catRM}

\dictentry{RP-1}{A refined kerosene used as rocket fuel that offers high density, ease of handling at ambient temperature, and good combustion characteristics when burned with liquid oxygen. RP-1 is less prone to coking than standard kerosene and provides excellent performance for sea-level engines where its high density reduces tank volume and structural weight. The F-1 engine on the Saturn V and the Merlin engine on the Falcon 9 both burn RP-1 with liquid oxygen. RP-1 formulations are carefully controlled to minimize soot formation and coking that could foul injector faces and cooling channels during extended engine operation.}{catRM}

\dictentry{S-400}{A Russian long-range air defense system, designated NATO reporting name Growler, capable of engaging aircraft, cruise missiles, and ballistic targets at ranges up to 400 kilometers. The S-400 uses multiple interceptor types optimized for different target altitudes and ranges, providing layered coverage within its engagement envelope. It features advanced phased-array radar and electronic counter-countermeasures that make it one of the most capable air defense systems in the world. The S-400 can simultaneously track hundreds of targets and engage dozens of them, providing robust defense against saturation attacks.}{catRM}

\dictentry{S-500}{A Russian next-generation air and missile defense system designed to engage hypersonic weapons, low-orbit satellites, and intercontinental ballistic missile warheads at ranges exceeding 600 kilometers. The S-500 features a new radar with enhanced detection capabilities and interceptors designed for exoatmospheric engagement of space-based threats. It is intended to supplement and eventually replace the S-400 as Russia's premier air defense system. The S-500 represents a significant advancement in capability, targeting the emerging threats of hypersonic glide vehicles and low-orbit space weapons.}{catRM}

\dictentry{Safe and Arm Device}{A safety device preventing accidental ignition by physically isolating the ignition train and ensuring that electrical or pyrotechnic energy cannot reach the initiator until the device is intentionally armed. Safe and arm devices use mechanical switches, thermal barriers, or rotary positioners to maintain a safe condition during handling, storage, and transport. They are mandatory on all military munitions and rocket motors to prevent inadvertent detonation. The device must provide a clear, verifiable safe condition that can be confirmed by visual inspection and a positive arm transition that requires deliberate action.}{catRM}

\dictentry{SAM}{A surface-to-air missile launched from ground-based platforms to engage aircraft, helicopters, cruise missiles, and ballistic targets. SAM systems range from man-portable shoulder-fired weapons like the Stinger to large long-range systems like the S-400 and Patriot. They use radar, infrared, or command guidance to steer toward the target and employ blast-fragmentation or hit-to-kill warheads for destruction. Modern SAM systems integrate multiple interceptor types and engagement modes to provide layered defense against diverse aerial threats across a wide range of altitudes and ranges.}{catRM}

\dictentry{Satellite Warning System}{A satellite system detecting missile launches using space-based infrared sensors that scan the Earth's surface for the intense heat signature of rocket motor exhaust plumes. Systems such as the US Space Based Infrared System provide global, persistent surveillance with detection times measured in seconds from boost phase ignition. The early warning data is relayed to ground stations for trajectory computation and dissemination to national command authorities and missile defense interceptors. Satellite warning systems provide the earliest possible detection of missile launches, maximizing the available warning time for defensive response.}{catRM}

\dictentry{Scene Matching}{Guidance matching imagery data stored in the missile's memory with real-time sensor observations of the terrain or target area below. The onboard computer compares pre-loaded reference images with actual camera or radar imagery to compute position corrections and steer toward the aim point. Scene matching provides high terminal accuracy independent of GPS and is used in cruise missiles and precision-guided munitions for area correlation navigation. The accuracy of scene matching depends on the quality of reference imagery, sensor resolution, and the distinctive features of the terrain being traversed.}{catRM}

\dictentry{Scramjet Missile}{A missile using a scramjet engine that maintains supersonic combustion of fuel in the airstream at hypersonic speeds, eliminating the need for onboard oxidizer and dramatically increasing range potential compared to rocket-powered missiles. Scramjets require speeds above Mach 5 to function and must be boosted to operating speed by a rocket or turbojet first stage. Several nations are developing scramjet missiles as a new class of hypersonic strike weapons with extended range and high speed. The sustained hypersonic flight creates extreme thermal environments that demand advanced materials and cooling strategies for the engine inlet, combustor, and nozzle.}{catRM}

\dictentry{Screaming}{High-frequency combustion instability characterized by pressure oscillations above 1000 Hz that can cause extremely rapid thermal and mechanical loading of engine components. Screaming is the most destructive form of combustion instability, capable of destroying an engine in milliseconds through burnthrough of chamber walls and injector faces. It is typically driven by tangential or radial acoustic modes and requires immediate design correction through injector pattern changes, baffles, or resonators. Screaming instability is notoriously difficult to predict and often first appears unexpectedly during hot-fire testing at full power conditions.}{catRM}

\dictentry{SECO}{Second engine cutoff, the moment when the second stage completes its planned burn and shuts down, typically after achieving the target orbit or trajectory. SECO timing is critical because overshooting or undershooting the target velocity directly affects the achieved orbit and mission success. Following SECO, the upper stage typically performs separation of the payload or prepares for a restart burn if the mission requires multiple burns. The precise determination of SECO time requires accurate real-time velocity measurement and comparison against the planned orbital parameters.}{catRM}

\dictentry{Second Stage}{The stage igniting after first stage separation, optimized for performance in near-vacuum conditions with a higher expansion ratio nozzle and lower thrust-to-weight requirement. The second stage takes over propulsion once the vehicle has cleared the dense atmosphere and accelerated beyond the capability of the first stage. It typically carries its own guidance system and may restart for orbit insertion or trajectory correction maneuvers. Second stage engines are optimized for vacuum specific impulse rather than sea-level thrust, using high-expansion-ratio nozzles that would be inefficient at lower altitudes.}{catRM}

\dictentry{Secondary Injection}{Injection of secondary fluid for thrust vectoring, where a reactive or inert fluid is injected into the nozzle divergent section to create a side force through localized pressure disturbances. This technique steers the thrust vector without gimbaling the engine and can be applied to both liquid and solid rocket motors. Secondary injection is simpler mechanically than gimbal systems but produces less control authority and some performance loss. The injection port location, fluid type, and flow rate must be carefully optimized to achieve the required control authority with acceptable thrust loss.}{catRM}

\dictentry{Semi-Active Radar Homing}{Homing using reflected radar from an external illuminator where the launch platform's radar illuminates the target and the missile's seeker detects the reflected energy to guide toward it. The missile does not carry its own transmitter, making it lighter, cheaper, and more covert than active radar homing missiles. The tradeoff is that the launch platform must maintain radar lock on the target until intercept, limiting its ability to engage other threats simultaneously. Semi-active radar homing remains widely used for its balance of performance, cost, and platform burden across many missile types.}{catRM}

\dictentry{Separation Motor}{A small motor separating stages by firing to push the spent stage away from the continuing upper stage, ensuring clean mechanical separation without recontact. Separation motors provide a positive axial impulse that opens the distance between stages quickly and predictably. They are typically solid-propellant devices mounted on the interstage structure and ignited by the flight control system at the commanded separation time. Separation motor thrust must be sufficient to overcome any residual aerodynamic or gravitational forces that could cause recontact between the separating stages.}{catRM}

\dictentry{Shaped Charge}{A shaped explosive for armor penetration that uses a metal-lined conical cavity to focus the detonation energy into a high-velocity metal jet capable of penetrating thick armor plate. The detonation wave collapses the liner inward, creating a hypersonic jet of copper or tungsten that erodes through armor by hydrodynamic impact. Shaped-charge warheads are the primary killing mechanism of anti-tank guided missiles and many other anti-armor weapons. The liner material, cone angle, and standoff distance are critical design parameters that determine the penetration depth and effectiveness against different armor types.}{catRM}

\dictentry{Shear Injector}{An injector using shear for atomization where fuel and oxidizer streams flow at high relative velocity past each other, and viscous shear forces break the liquid into fine droplets. Shear atomization is effective when one or both propellants have high viscosity or surface tension that resists breakup by impingement alone. This injector type is used in some pressure-fed engines where simplicity and reliable atomization across a range of flow rates are important. The shear injector design provides good atomization without the precise alignment requirements of impinging jet injectors.}{catRM}

\dictentry{Showerhead Injector}{A simple injector with many small holes arranged in a flat pattern across the injector face, producing individual streams of propellant that mix and combust in the chamber volume. Showerhead injectors are easy to manufacture and provide uniform propellant distribution but may have lower atomization quality compared to impinging or coaxial designs. They are commonly used in small thrusters, gas generators, and early liquid rocket engines where simplicity is valued. The showerhead pattern must be carefully designed to prevent coherent acoustic modes that could couple with combustion and drive instability.}{catRM}

\dictentry{Side Load}{Lateral loads from flow separation in a nozzle caused by asymmetric detachment of the exhaust flow from the nozzle wall, typically during sea-level operation of overexpanded nozzles. Side loads can produce moments that stress the nozzle, gimbal bearings, and vehicle structure, and in extreme cases cause nozzle failure. Nozzle design must account for side load conditions, and testing includes characterization of separation onset and lateral force magnitude. Side loads are most severe during engine start and shutdown transients when the chamber pressure passes through the range most susceptible to asymmetric flow separation.}{catRM}

\dictentry{Silo}{An underground launch facility for missiles consisting of a hardened concrete and steel shaft equipped with environmental control, propellant loading, and launch sequencing systems. Silo-based ICBMs can be launched within minutes of receiving the launch order, providing rapid response capability as a key element of nuclear deterrence. Modern silos are designed to withstand near-miss nuclear blasts and incorporate shock isolation systems to protect the missile from ground shock. Silo launch systems include exhaust management provisions to safely redirect the missile exhaust gases during the critical initial launch phase within the confined shaft.}{catRM}

\dictentry{SLBM}{A submarine-launched ballistic missile fired from submerged submarines to deliver nuclear or conventional warheads against land or sea targets. SLBMs provide the most survivable leg of the nuclear triad because submarines can patrol undetected in the deep ocean, making a retaliatory strike possible even after a first strike. Current examples include the US Trident II D5 and the Russian Bulava, both capable of carrying multiple independently targetable warheads. SLBM design must accommodate the unique challenges of submerged launch, including water-exit dynamics, gas bubble management, and missile control during the transition from underwater to atmospheric flight.}{catRM}

\dictentry{Slotted Tube Grain}{A tubular grain with slots for burning surface control where axial slots in a cylindrical perforation create additional burning surfaces and control the thrust-time profile. The slot geometry provides progressive burning characteristics by increasing the burning surface area as the slots widen during combustion. Slotted tube grains are widely used in tactical missile motors where specific thrust profiles must be achieved within compact motor dimensions. The slot width, length, and number are optimized during motor design to produce the required thrust-time history while maintaining grain structural integrity.}{catRM}

\dictentry{Solenoid Valve}{An electrically actuated valve that uses an electromagnetic coil to move a plunger that opens or closes the flow path, providing rapid and precise control of propellant or gas flow. Solenoid valves offer fast response times, typically under 50 milliseconds, and can be controlled directly by digital flight computers without hydraulic or pneumatic actuators. They are widely used for propellant sequencing, purge control, and small thruster valve actuation in rocket propulsion systems. Solenoid valve power consumption and thermal management are critical design considerations for battery-powered spacecraft applications.}{catRM}

\dictentry{Solid Rocket Booster}{A solid rocket used as a strap-on booster attached to the first stage of a launch vehicle to provide additional thrust during the initial boost phase. SRBs are jettisoned after burnout, typically 2 minutes into flight, and are recovered by parachute for refurbishment and reuse in some systems. The Space Shuttle and SLS use large reusable SRBs, while many expendable launch vehicles use smaller one-piece solid boosters for payload augmentation. SRB thrust profile, burn duration, and separation characteristics must be carefully matched to the core vehicle requirements for optimal mission performance.}{catRM}

\dictentry{Solid Rocket Motor}{A rocket using solid chemical propellant cast in a case that provides thrust when the grain is ignited and burns from the exposed surfaces inward. Solid motors offer simplicity, high reliability, and immediate launch readiness because the propellant is pre-loaded and requires no fueling sequence before firing. They are used extensively in tactical missiles, launch vehicle boosters, and strategic ICBMs where rapid response and operational simplicity are paramount. Solid motor performance is determined during the design and casting phase, with no ability to adjust thrust or burn duration after manufacture.}{catRM}

\dictentry{Sounding Rocket}{A suborbital rocket for research that carries scientific instruments to the upper atmosphere and near-space environment to conduct experiments in microgravity, atmospheric chemistry, and solar physics. Sounding rockets follow a parabolic trajectory, providing several minutes of microgravity at the peak of their flight before returning to Earth by parachute. They are a cost-effective alternative to satellites for short-duration experiments and technology demonstrations. Sounding rocket programs provide valuable flight heritage for new technologies and serve as an accessible entry point for student and university research in space science.}{catRM}

\dictentry{Spark Igniter}{An igniter using an electric spark generated by a spark plug or torch igniter to initiate combustion of propellants in the chamber. Spark ignition is clean, repeatable, and allows multiple restarts, making it ideal for upper-stage and orbital maneuvering engines. The RL10 and Vinci engines use spark ignition systems that provide reliable startup throughout the mission without the consumables required by pyrotechnic igniters. Spark igniter systems require a reliable electrical power source and ignition circuit that can deliver consistent spark energy across the full range of operating conditions and altitudes.}{catRM}

\dictentry{Specific Impulse}{A measure of rocket propellant efficiency defined as the thrust produced per unit weight flow of propellant, expressed in seconds. Higher specific impulse means more thrust per unit of propellant consumed, directly translating to greater delta-v and payload capacity for a given vehicle mass. Typical values range from about 200 seconds for solid motors to over 450 seconds for liquid oxygen and liquid hydrogen engines. Specific impulse is the single most important performance parameter for comparing different propulsion systems and propellant combinations.}{catRM}

\dictentry{Specific Impulse Efficiency}{The achieved Isp as a fraction of theoretical, comparing the actual specific impulse measured during hot-fire testing to the ideal value calculated from thermodynamic equilibrium for the propellant combination. Isp efficiency reflects losses from incomplete combustion, nozzle flow divergence, boundary layer effects, and chemical kinetics. High-performance engines typically achieve 90 to 95 percent of theoretical Isp efficiency. The gap between theoretical and achieved Isp is a primary focus of engine development, with incremental gains achieved through injector optimization, chamber design improvements, and combustion process refinements.}{catRM}

\dictentry{Spin Motor}{A motor spinning a stage for stability that applies a tangential impulse to impart a controlled roll rate, providing gyroscopic stability to the missile or rocket during its unguided flight phase. Spin stabilization eliminates the need for fin stabilization or active attitude control, simplifying the motor design and reducing cost. Spin rates of several hundred revolutions per minute are typical for spin-stabilized rockets and some artillery projectiles. The spin motor must deliver a precise and repeatable impulse to achieve the target spin rate without exceeding the structural limits of the spinning component.}{catRM}

\dictentry{Spin-Stabilized Rocket}{A rocket stabilized by rotating around its longitudinal axis that uses gyroscopic rigidity to resist disturbing torques and maintain its pointing direction during flight. The spin rate must be carefully controlled to avoid resonance with structural modes and to ensure the guidance system can measure and correct for any deviations. Spin stabilization is simple and effective for short-range rockets and some missile types where active control is impractical. The spin rate must be high enough to provide adequate stability but low enough to avoid centrifugal loads that could affect guidance sensor accuracy.}{catRM}

\dictentry{SRBM}{A short-range ballistic missile with a range under 1000 kilometers that follows a ballistic trajectory to deliver conventional or nuclear warheads against theater targets. SRBMs are the most numerous category of ballistic missiles worldwide and are deployed by many nations on mobile launchers, in silos, and on naval vessels. Their short flight time of just a few minutes makes them particularly difficult to detect and defend against once launched. The proliferation of SRBMs represents a significant regional security challenge due to their affordability, accuracy, and rapid deployment capability.}{catRM}

\dictentry{Stage Separation}{The process of discarding an expended stage when the lower stage has depleted its propellant and is no longer contributing to vehicle acceleration. Separation is accomplished by pyrotechnic bolts, explosive charges, or pneumatic pushers that sever the structural connection and push the stages apart. Clean and reliable stage separation is critical because any recontact between stages can cause catastrophic vehicle destruction. Separation systems must provide symmetric, controlled push-off forces to ensure clean separation without inducing tumbling or tip-off rates that could affect the subsequent stage ignition.}{catRM}

\dictentry{Staged Combustion Cycle}{A cycle burning preburner exhaust in the main chamber where propellant is partially reacted in one or more preburners to drive the turbopump, and then the turbine exhaust is routed into the main chamber for complete combustion. This approach extracts maximum energy from the propellant by avoiding the dump loss of a gas generator cycle. Staged combustion engines achieve the highest specific impulse and chamber pressures of any chemical rocket cycle, as demonstrated by the SSME and RD-180. The high chamber pressures enabled by this cycle require robust turbopump and combustion chamber designs that can withstand extreme operating conditions.}{catRM}

\dictentry{Star Grain}{A grain with a star-shaped internal bore that provides a progressive-burning characteristic by creating a large initial surface area in the star points that decreases as the bore expands. Star grains are one of the most common grain configurations in tactical missile motors, offering excellent thrust tailoring in a compact geometry. The star pattern dimensions and point count are optimized during motor design to achieve the desired thrust-time history. Star grain design provides a good balance between thrust profile flexibility and manufacturing simplicity for high-production tactical motor applications.}{catRM}

\dictentry{Strand Burn Rate}{Burn rate measured on a small sample of propellant in a strand burner under controlled pressure and temperature conditions to characterize the material's burning behavior. The strand test provides fundamental burn rate data as a function of pressure and initial temperature, which is used to predict full-scale motor performance. This data feeds into the burn rate law of the form r equals a times P to the power n, where the pressure exponent n is critical for motor stability assessment. Strand burn rate data is essential for validating propellant formulations and ensuring consistency across production lots.}{catRM}

\dictentry{Strap-On Booster}{An auxiliary booster attached to the first stage that provides additional thrust during lift-off and early flight before being jettisoned. Strap-on boosters can be solid or liquid and are added in varying numbers to tailor the launch vehicle's thrust and payload capability for specific missions. Examples include the SRBs on the Space Shuttle and SLS, and the Castor solid strap-on boosters used on Delta and Atlas vehicles. The timing and sequence of strap-on booster ignition and separation must be carefully coordinated with the core vehicle to optimize trajectory performance.}{catRM}

\dictentry{Suction Line}{A line on the inlet side of the pump that carries low-pressure propellant from the tank to the turbopump inlet. The suction line must be designed to prevent cavitation by maintaining adequate inlet pressure and avoiding flow separation or vapor formation at the pump inducer. Its diameter is typically larger than the discharge line to keep flow velocities low and minimize pressure drop. Suction line design must account for propellant temperature, tank head pressure, and acceleration environments to ensure bubble-free delivery to the pump inlet under all operating conditions.}{catRM}

\dictentry{T-0}{The moment of engine ignition or lift-off, the designated zero time from which all events in the launch countdown and flight timeline are referenced. T-0 is precisely synchronized across all ground systems, tracking networks, and onboard avionics to ensure coordinated sequencing of engine start, hold-down release, and umbilical retraction. The definition of T-0 varies between ignition of the main engines and actual release from the pad depending on the launch vehicle design. Accurate T-0 timing is essential for trajectory reconstruction, performance analysis, and range safety monitoring during the launch campaign.}{catRM}

\dictentry{Tail Fins}{Aerodynamic surfaces on the rear of a rocket for stability that provide restoring forces when the rocket deviates from its velocity vector, keeping the nose pointed into the relative wind. Tail fins shift the center of pressure aft of the center of gravity, creating a natural static stability margin. Fixed tail fins are the simplest and most reliable form of flight stabilization, widely used on unguided rockets, missiles, and artillery projectiles. The fin planform, sweep angle, and aspect ratio are optimized to provide sufficient stability while minimizing drag and weight penalties during the ascent trajectory.}{catRM}

\dictentry{Tank Pressurization}{Pressurizing tanks to feed propellant to the engine by introducing high-pressure gas above the liquid propellant to maintain the required inlet pressure at the pump or injector. Pressurization methods include helium stored in composite overwrapped bottles, autogenous pressurization using vaporized propellant, and electric pump-fed systems that eliminate the need for tank pressure. The pressurization system must maintain stable tank pressure throughout the burn despite changing propellant levels and thermal conditions. Tank pressurization transients during engine start and shutdown must be carefully controlled to avoid propellant feed disturbances that could affect combustion stability.}{catRM}

\dictentry{Tap-Off Cycle}{A cycle tapping combustion gas to drive the turbine by extracting a portion of hot combustion products directly from the main chamber at a point downstream of the injector. The tapped gas is at lower temperature than a dedicated preburner output, which reduces turbine thermal loading. The tap-off cycle was used on the J-2 engine of the Saturn V and is conceptually similar to the gas generator cycle but draws gas from the main chamber. This cycle offers a compromise between the simplicity of gas generator and the performance of staged combustion cycles.}{catRM}

\dictentry{Target Drone}{An unmanned target used for missile testing that simulates the radar, infrared, and flight characteristics of real threats such as aircraft, cruise missiles, or ballistic reentry vehicles. Target drones can be recovered by parachute for reuse or expended during live-fire exercises. They provide a realistic training environment for air defense crews and are also used during missile development testing to validate guidance and warhead performance. Modern target drones can replicate complex evasive maneuvers and deploy countermeasures to provide realistic engagement scenarios for missile system evaluation.}{catRM}

\dictentry{Terminal Guidance}{The final phase of missile guidance where the onboard seeker acquires the target and generates steering commands to achieve a direct hit or optimal detonation geometry. Terminal guidance is the most accuracy-critical phase, as the missile must correct any residual midcourse errors within the short time remaining before impact. Different seekers such as radar, infrared, laser, or GPS provide the sensing capability for various target types and engagement scenarios. The terminal guidance algorithm must process sensor data and generate control commands at high rates to achieve the required miss distance for effective warhead detonation.}{catRM}

\dictentry{Terminal Phase}{The final engagement phase of an interceptor when the kill vehicle separates from its booster and uses onboard sensors and divert thrusters to home in on and collide with the target. The terminal phase lasts only seconds and requires extremely fast sensor processing and maneuvering capability because the closing velocity between interceptor and target may exceed 10 kilometers per second. Success in the terminal phase depends on the accuracy of midcourse guidance that placed the interceptor on a near-intercept trajectory. The kill vehicle must autonomously acquire and track the target with sufficient accuracy to generate divert commands for the final homing maneuver.}{catRM}

\dictentry{Terrain Contour Matching}{Guidance matching terrain data by comparing barometric or radar altimeter measurements of the ground profile beneath the missile with pre-stored digital terrain elevation maps. The onboard computer correlates the measured terrain cross-section with the stored map to determine position errors and generate steering corrections. Terrain contour matching, used in cruise missiles like the Tomahawk, provides GPS-independent navigation accuracy of tens of meters. TCAM accuracy depends on the distinctiveness of the terrain features along the flight path, with flat or repetitive terrain providing less discriminating capability.}{catRM}

\dictentry{THAAD}{Terminal High Altitude Area Defense, a US Army missile defense system designed to intercept short-, medium-, and intermediate-range ballistic missiles in their terminal phase at altitudes between 40 and 150 kilometers. THAAD uses a hit-to-kill interceptor guided by the AN/TPY-2 X-band radar, providing a critical layer in the US ballistic missile defense architecture. The system can be deployed rapidly by air transport and is designed to protect military forces and population centers from theater ballistic missile threats. THAAD's upper-tier engagement capability complements lower-tier systems like Patriot by intercepting threats at higher altitudes before they descend into densely populated areas.}{catRM}

\dictentry{Throat}{The narrowest section of a rocket nozzle where the flow reaches sonic velocity (Mach 1) and the mass flow rate is limited by the throat area and upstream conditions. The throat area is the most critical dimension in nozzle design because it determines the chamber pressure for a given propellant mass flow rate and the overall thrust level. Throat erosion during operation increases the effective area, reducing chamber pressure and degrading performance over the burn. The throat contour must be precisely manufactured to ensure uniform flow acceleration and avoid local hot spots or flow disturbances.}{catRM}

\dictentry{Thrust}{The force propelling a rocket forward, generated by expelling mass at high velocity through the nozzle according to Newton's third law. Thrust equals the product of mass flow rate and exhaust velocity plus the pressure difference at the nozzle exit multiplied by exit area. The total impulse, which is thrust integrated over burn time, determines the total velocity change the rocket can deliver to its payload. Thrust vector direction and magnitude must be precisely controlled throughout the burn to follow the planned trajectory and achieve the target orbit or impact point.}{catRM}

\dictentry{Thrust Chamber}{The assembly including combustion chamber and nozzle that converts propellant chemical energy into thrust through combustion and expansion. The thrust chamber is the core component of a liquid rocket engine, integrating the injector, chamber, throat, and nozzle into a single functional unit. Its design must simultaneously satisfy requirements for performance, cooling, structural integrity, combustion stability, and manufacturing feasibility. Thrust chamber development involves extensive analytical modeling, cold-flow testing, and hot-fire validation to demonstrate that all performance and durability requirements are met.}{catRM}

\dictentry{Thrust Chamber Liner}{A liner protecting the thrust chamber wall from the direct impingement of hot combustion gases, often used in regeneratively cooled engines as an additional thermal barrier between the coolant channels and the combustion environment. The liner is typically made from a copper alloy or nickel superalloy that is brazed or bonded to the coolant jacket structure. It must withstand high thermal gradients, erosion, and cyclic loading throughout the engine's operational life. The liner thickness and material selection are optimized to provide adequate thermal protection while minimizing the thermal resistance between the coolant and the combustion gas environment.}{catRM}

\dictentry{Thrust Coefficient}{A dimensionless measure of nozzle performance defined as the ratio of actual thrust to the thrust that would be produced if the chamber pressure acted on the throat area. Thrust coefficient depends on the nozzle expansion ratio, the ratio of specific heats of the combustion products, and the degree of nozzle expansion relative to ambient pressure. Higher thrust coefficients indicate more efficient conversion of chamber pressure into useful thrust. The thrust coefficient is a key parameter in the thrust equation and is used to compare nozzle performance across different engine designs and operating conditions.}{catRM}

\dictentry{Thrust Profile}{The variation of thrust over burn time, showing how the engine's force output changes from ignition to burnout. Thrust profiles are shaped by the grain geometry in solid motors and by the throttle schedule in liquid engines, and can be progressive, neutral, or regressive. The thrust profile determines the vehicle's acceleration history, structural loads, and trajectory, and is a primary output of the motor design process. Matching the achieved thrust profile to the design specification is a critical qualification criterion for both solid and liquid propulsion systems.}{catRM}

\dictentry{Thrust Termination}{A system to shut down thrust in a solid rocket motor, typically by venting chamber pressure through ports in the forward end or side of the case to drop pressure below the sustaining level. Thrust termination is used for range control in tactical missiles and for precise cutoff velocity in launch vehicle upper stages. Once initiated, thrust termination is irreversible because solid propellant cannot be extinguished by simply stopping a feed system. The termination port design must ensure rapid and complete pressure decay to prevent thrust residual that could affect the trajectory after the planned burnout point.}{catRM}

\dictentry{Thrust Vector Control}{Steering a rocket by redirecting thrust using gimbaled nozzles, secondary injection, jet vanes, or differential thrust from multiple engines. TVC provides the control forces and moments needed to follow the planned trajectory and correct for wind disturbances and asymmetric thrust conditions. It is the primary means of steering during powered flight for both liquid and solid rocket-powered vehicles. TVC system response time and control authority must be sufficient to handle all expected disturbance environments while maintaining trajectory accuracy within mission requirements.}{catRM}

\dictentry{Thrust-to-Weight Ratio}{The ratio of engine thrust to vehicle weight at lift-off, which must exceed 1.0 for the vehicle to leave the ground and is typically between 1.2 and 1.5 for launch vehicles. Higher thrust-to-weight ratios produce faster acceleration through the lower atmosphere, reducing gravity and drag losses, but require more structural strength to handle the increased loads. Thrust-to-weight ratio changes during flight as propellant is consumed and the vehicle becomes lighter. The initial thrust-to-weight ratio at lift-off is a critical design parameter that influences the entire ascent trajectory and structural design requirements.}{catRM}

\dictentry{Time Fuse}{A fuse detonating after a set time that uses a mechanical clockwork, pyrotechnic delay element, or electronic timer to initiate the warhead at a predetermined time after launch or arming. Time fuzes are used for airburst effects, fuzing of proximity-capable munitions, and self-destruct functions to prevent unexploded ordnance from littering the battlefield. Modern electronic time fuzes offer microsecond precision and can be programmed before launch based on the expected time of flight. Time fuse accuracy directly affects the height of burst and the lethality coverage area against the intended target type.}{catRM}

\dictentry{Total Impulse}{The total thrust integrated over burn time, representing the total impulse delivered by the motor and measured in Newton-seconds. Total impulse determines the total delta-v capability of the rocket and is the primary metric for sizing solid rocket motors for a given mission. Motors are classified by total impulse ranges such as miniature, small, medium, and large, which correspond to different tactical and strategic applications. Total impulse is the product of average thrust and burn time, and both parameters must be optimized together to meet the mission delta-v requirements.}{catRM}

\dictentry{Tracking Radar}{Radar following a target by continuously measuring its range, bearing, elevation, and velocity to provide fire control data for weapon engagement. Tracking radars use narrow beams and closed-loop servo mechanisms to maintain lock on the target with high accuracy. They provide the targeting data for semi-active, command-guided, and trajectory-predicting missile engagement modes. Tracking radar accuracy and update rate directly determine the quality of the fire control solution and the effectiveness of the missile engagement against maneuvering targets.}{catRM}

\dictentry{Trajectory}{The path followed by a missile or rocket through space and the atmosphere, determined by initial conditions, propulsion, guidance commands, gravitational forces, and aerodynamic effects. Trajectory design balances competing requirements such as range, flight time, thermal loads, structural loads, and defense penetration against time of flight constraints and propellant budgets. The chosen trajectory type, whether ballistic, lofted, depressed, or lifting, has profound implications for weapon effectiveness and missile defense. Trajectory optimization involves iterative analysis of the propulsion, aerodynamic, and gravitational effects throughout the entire flight profile.}{catRM}

\dictentry{Tsiolkovsky Equation}{The rocket equation delta-v equals ve times the natural log of initial mass over final mass, which relates the achievable velocity change to the exhaust velocity and the mass ratio of the vehicle. This fundamental relationship shows that delta-v increases only logarithmically with mass ratio, making it very difficult to achieve large velocity changes with a single stage. The Tsiolkovsky equation is the foundational design tool for sizing rocket stages and determining payload capability. Understanding the Tsiolkovsky equation is essential for appreciating why staging is necessary for orbital flight and why propellant mass fraction is such a critical design parameter.}{catRM}

\dictentry{Turbopump}{A turbine-driven pump feeding propellant to the engine, consisting of one or more centrifugal or axial pumps driven by a gas turbine to deliver propellant at the required flow rate and pressure. Turbopumps are among the most challenging components in a rocket engine because they must handle cryogenic or reactive fluids at extremely high rotational speeds while maintaining reliability. They enable high chamber pressures and thrust levels that are not achievable with pressure-fed systems. Turbopump development represents one of the most technically demanding aspects of liquid rocket engine design, requiring expertise in fluid dynamics, materials science, and high-speed rotating machinery.}{catRM}

\dictentry{UDMH}{Unsymmetrical dimethylhydrazine, a hypergolic fuel with the chemical formula C2H8N2 that is commonly used with nitrogen tetroxide in storable bipropellant rocket engines. UDMH offers good performance, liquid range from minus 57 to 63 degrees Celsius, and long-term storability, making it suitable for ballistic missile and spacecraft propulsion. Despite its effectiveness, UDMH is highly toxic and carcinogenic, requiring strict handling protocols and limiting its use in modern systems. UDMH is gradually being replaced by less toxic alternatives in new propulsion system designs where environmental and health concerns are prioritized.}{catRM}

\dictentry{Ullage Motor}{A small motor settling propellant before main engine ignition by producing a small forward acceleration that pushes liquid propellant to the tank outlet in microgravity or during coast phases. Ullage motors ensure that the main engine draws gas-free liquid propellant upon startup, preventing pump cavitation and combustion instability. They are typically small solid or cold-gas thrusters that fire for several seconds before the main engine ignition sequence. The ullage impulse must be sufficient to overcome surface tension and residual acceleration effects to ensure complete liquid coverage of the tank outlet.}{catRM}

\dictentry{Underexpansion}{Nozzle exit pressure above ambient, occurring when a nozzle designed for high-altitude or vacuum operation is fired at lower altitude where atmospheric pressure is lower than the exhaust exit pressure. Underexpanded nozzles waste potential thrust because the exhaust could be expanded further, but the condition is acceptable during the initial ground-level phase of launch if the engine is optimized for vacuum performance. Extended nozzle bells or plug nozzles can recover some of the underexpansion loss at low altitude. The degree of underexpansion decreases as altitude increases and ambient pressure drops toward the nozzle exit pressure.}{catRM}

\dictentry{Upper Stage}{The final stage for orbit insertion, optimized for vacuum operation with high-expansion-ratio nozzles and restart capability for precise orbit circularization, plane changes, and multiple payload deployment. Upper stages operate in near-vacuum conditions where specific impulse is maximized and must function reliably after a coast period following stage separation. Examples include the Centaur, Fregat, and Delta cryogenic upper stages that provide high-performance orbit insertion capability. Upper stage design must balance the performance benefits of high-expansion-ratio nozzles against the practical constraints of stage length, mass, and structural loads during the coast and restart phases.}{catRM}

\dictentry{Vent Valve}{A valve for tank venting that releases excess pressure from propellant or pressurant tanks during ground operations, ascent, and thermal conditioning to prevent overpressurization. Vent valves are typically spring-loaded or pyrotechnically operated and are sized to handle the maximum expected pressurization rate from thermal expansion or propellant loading. They are a critical safety component that prevents tank rupture while maintaining the required pressurization during flight. Vent valve sizing must account for the worst-case pressurization scenarios including solar heating, propellant transfer, and pressurant gas thermal expansion.}{catRM}

\dictentry{Vernier Engine}{A small engine for fine trajectory adjustments that provides low thrust for precise velocity corrections during orbit insertion, station-keeping, and rendezvous maneuvers. Vernier engines are typically restartable hypergolic or monopropellant thrusters that complement the main engine's coarse thrust with fine control authority. They are essential for achieving the precise orbit parameters required for satellite constellation deployment and planetary missions. Vernier engine performance is characterized by very fine thrust resolution and high impulse bit accuracy to enable the precise velocity adjustments needed for final orbit circularization.}{catRM}

\dictentry{Wagon Wheel Grain}{A grain with a wagon wheel cross-section featuring a central bore with radiating spoke-like perforations that create a large and complex burning surface. This geometry provides progressive-neutral burning characteristics by balancing the increasing surface of the widening spokes against the decreasing outer surface. Wagon wheel grains are used in tactical missile motors where specific thrust shaping is required within tight dimensional constraints. The spoke width, length, and angular spacing are optimized during motor design to produce the required thrust-time history while maintaining adequate grain structural integrity.}{catRM}

\dictentry{Warhead}{The explosive or payload of a missile that inflicts damage on the target through blast overpressure, fragmentation, shaped-charge jet, kinetic energy, or nuclear effects. Warhead design must be matched to the target type, fuze mode, and engagement geometry to maximize lethality. Common warhead types include blast-fragmentation, continuous rod, shaped-charge, and hit-to-kill kinetic kill vehicles, each optimized for a different class of target. Warhead effectiveness is validated through extensive testing against representative targets to verify that the designed lethal mechanisms perform as predicted in realistic engagement conditions.}{catRM}

\dictentry{Wet Mass}{The mass of the rocket with propellant, representing the total launch weight including structure, engines, avionics, payload, and all propellant. Wet mass is the initial condition for the rocket equation and determines the thrust-to-weight ratio at lift-off. Minimizing wet mass for a given payload and delta-v requirement is the fundamental objective of launch vehicle structural optimization. Wet mass determines the gravitational and inertial loads experienced by the vehicle structure throughout the flight and directly influences the sizing of structural members, tank walls, and engine mounts.}{catRM}

\dictentry{Wire Guidance}{Guidance via wires spooled from the missile that connects the weapon to the launch platform through thin copper or fiber-optic wires unspooled during flight. Steering commands are transmitted over the wire from the operator or fire control system, providing precise terminal guidance against moving or fortified targets. Wire guidance is highly resistant to electronic jamming and is used in anti-tank missiles like the TOW and ERYX where reliability and precision are paramount. The wire link must maintain integrity throughout the missile's flight while the wire unspools at high speed without tangling or breaking.}{catRM}

\dictentry{World Wide Web}{Not a rocket term. See: Web, World Wide.}{catRM}

\dictentry{WSB Trajectory}{Weak Stability Boundary trajectory, a low-energy transfer path that exploits the chaotic dynamics near Lagrange points to move between orbits with minimal propellant. WSB trajectories use the gravitational interplay between three bodies to achieve transfers that would be impossible with conventional propulsion. The Hiten and Genesis missions demonstrated WSB transfers to the Moon with significantly reduced delta-v compared to Hohmann transfers.}{catRM}

\dictentry{X-Ray Pulsar Navigation}{An autonomous navigation technique for deep-space spacecraft that uses the precisely timed X-ray pulses from known pulsars as natural beacons. By measuring the arrival times of pulses from multiple pulsars, a spacecraft can determine its position and velocity without ground-based tracking. X-ray pulsar navigation is being developed for future interplanetary missions where ground tracking is impractical due to light-time delays and limited DSN availability.}{catRM}

\dictentry{Xenon}{A heavy noble gas used as propellant in electric propulsion systems, particularly Hall-effect thrusters and ion engines. Xenon's high atomic mass, low ionization energy, and chemical inertness make it ideal for electrostatic acceleration, yielding high specific impulse with minimal erosion. It is stored as a compressed gas or supercritical fluid and is the standard propellant for commercial and scientific electric propulsion missions including Dawn, BepiColombo, and Starlink satellites.}{catRM}

\dictentry{X-Band}{A portion of the electromagnetic spectrum in the frequency range from approximately 8 to 12 gigahertz, used for high-data-rate spacecraft communication and radar applications. X-band offers a good compromise between the high data rates of Ka-band and the weather resilience of lower-frequency bands, making it widely used for deep-space links and military radar. X-band frequencies are also used for planetary radar mapping and spacecraft tracking.}{catRM}

\dictentry{Yaw Control}{Control of the missile's rotation about its vertical axis, achieved through thrust vector control, aerodynamic fins, or reaction control thrusters. Yaw control is essential for maintaining the desired heading during powered flight and for pointing the seeker or warhead during terminal guidance. It must coordinate with pitch and roll control to achieve the commanded attitude without coupling effects.}{catRM}

\dictentry{Yaw Plane}{The plane containing the missile's longitudinal and vertical axes, in which yaw rotation occurs. Yaw plane maneuvers change the missile's heading angle relative to the velocity vector and are used for azimuth corrections during midcourse and terminal guidance. The yaw plane is orthogonal to the pitch plane and together they define the full attitude space of the missile.}{catRM}

\dictentry{Ytterbium}{A rare earth element occasionally investigated as an alternative propellant for electric thrusters, particularly for field-emission electric propulsion (FEEP) where its low melting point and low ionization potential offer advantages. Ytterbium FEEP thrusters can operate at very low power levels with high specific impulse, suitable for precision formation flying and drag-free control.}{catRM}

\dictentry{Zarya}{The Functional Cargo Block (FGB) of the International Space Station, also known as Zarya (Dawn), launched in 1998 as the first ISS module. It provided initial propulsion, power, and storage for the early station assembly phase. Zarya was built by Russia under U.S. contract and represents a key milestone in international space cooperation.}{catRM}

\dictentry{Zenith}{The direction pointing directly upward from a location on Earth's surface, opposite to nadir. In rocket launches, the zenith direction is the initial ascent vector for vertical launch trajectories. Zenith angle is the angle between the zenith and the line of sight to a celestial object, used in ground-based tracking and launch azimuth calculations.}{catRM}

\dictentry{Zero-Lift Drag}{The drag force acting on a missile or rocket at zero angle of attack, comprising skin friction drag and pressure drag from the body shape. Zero-lift drag is the baseline drag that must be overcome at all flight conditions and is minimized through streamlined body design, smooth surface finishes, and area ruling. It is a primary factor in determining the vehicle's maximum speed and range.}{catRM}

\dictentry{Zinc-Air Battery}{A primary battery using zinc and atmospheric oxygen as reactants, offering high energy density for its weight. Zinc-air batteries have been used in some missile power supplies where long shelf life and high specific energy are required. The battery activates when air enters the cell, making it suitable for applications where the power source must remain inert until deployment.}{catRM}

\dictentry{ZND Model}{Zeldovich-von Neumann-Doering model, a one-dimensional theoretical model of detonation structure that describes the steady propagation of a detonation wave through a reactive medium. The model consists of a leading shock wave (von Neumann spike) followed by a reaction zone where chemical energy is released. The ZND model is fundamental to understanding detonation physics in solid and liquid propellants and in pulse detonation engines.}{catRM}

\dictentry{Zoning}{The practice of dividing a solid rocket motor grain into distinct sections with different propellant formulations or geometries to tailor the thrust-time profile. Zoning allows the motor designer to optimize performance for different flight phases, such as a high-thrust boost phase followed by a lower-thrust sustain phase. Each zone is cast or bonded separately and the interfaces must be carefully designed to prevent debonding or unintended ignition.}{catRM}

\dictentry{Zirconium}{A transition metal used in pyrotechnic initiators and as an additive in some solid propellants to increase flame temperature and combustion stability. Zirconium's high reactivity and exothermic oxidation make it effective in ignition systems and as a burn rate modifier. Zirconium-sponge pyrotechnics are used in NASA standard initiators and other space-qualified ignition devices.}{catRM}

